\documentclass[twocolumn,prx,aps,superscriptaddress,longbibliography]{revtex4-2}
\usepackage{amsmath,amssymb,amsthm}
\usepackage{graphicx}
\usepackage{booktabs}
\usepackage{hyperref}

\newtheorem{proposition}{Proposition}
\newtheorem{corollary}[proposition]{Corollary}
\newtheorem*{remark}{Remark}

% Claim-engine markers: zero-width in production.
\newcommand{\pclaim}[4]{}

\begin{document}

\title{Supplemental Material: Provable measurement design for quantum state tomography\\
from Hamming-scheme spectral theory}

\author{Andrew Craton}
\affiliation{Kootru Labs, Burlington, Wisconsin, USA}
\email{acraton@kootru.com}

\maketitle

This Supplemental Material provides extended derivations and analyses
supporting the main text.  Section~\ref{sm:summary} contains a summary table.
Sections~\ref{sec:qudit}--\ref{sec:extended-limitations} provide the
full technical details of extensions and robustness checks referenced
in the main text and Appendices~A--F.
Section~\ref{sm:open} lists open questions and Sec.~\ref{sm:data-code}
provides data and code availability information.

%% ============================================================
%% SUMMARY
%% ============================================================
\section{Summary}
\label{sm:summary}

\begin{table*}
\caption{Summary of supplemental analyses.  Columns indicate
whether the finding \emph{helps} the paper's claims, is
\emph{neutral}, or identifies a \emph{limitation}.}
\label{tab:sm-summary}
\begin{ruledtabular}
\begin{tabular}{p{0.20\textwidth}p{0.35\textwidth}lp{0.25\textwidth}}
Analysis & Key finding & Valence & Impact \\
\hline
Three-layer structure (main text, Appendix~C)
  & Spectral $\to$ lattice $\to$ combinatorial
  & Helps & Explains why $K^* = 5$ is not ad hoc \\
Circuit compression (main text, Appendix~D)
  & 29 circuits (not 137); 76.9\% coverage
  & Helps & 79\% shot reduction \\
Classical shadows (main text, Appendix~D)
  & $K^*$ = spectral front end for shadow tomography
  & Helps & Connects to established framework \\
RDM fidelities (main text, Appendix~E)
  & 1-RDM $F = 0.985$, 2-RDM $F = 0.771$
  & Helps & Strongest unreported practical result \\
Fisher/frame (main text, Appendix~F)
  & All Pauli subsets equivalent; weight-2 progressive addition $\Delta F = +0.204$
  & Helps & Standard criteria fail; weight-class completeness is the driver \\
MLE mechanism (main text, Appendix~F)
  & Positivity $+$ local completeness $\to$ convergence
  & Helps & Explains nonlinear advantage \\
Qudit extension (Sec.~\ref{sec:qudit})
  & $K^* = q^2$ ($q \ge 3$); $K = q^2{+}1$ for wt-1 sat.; 1-RDM $F = 0.972$ (sim)
  & Helps & Compositional $2$-qutrit patches validated in simulation \\
Dimensional proof (Sec.~\ref{sec:dim-proof})
  & Low-weight saturation at $K^* \geq 5$; cross-checked 3 reconstructors
  & Helps & Explains $\Delta(S{-}AR)$ drop at $n = 4$ \\
Ordering robustness (Sec.~\ref{sec:ordering})
  & Bimodal: 4/7 orderings at $F \approx 0.80$; default at median
  & Helps & Rules out cherry-picking; weight budget is the invariant \\
Near-optimality (Sec.~\ref{sec:optimality-test})
  & 90th pctl single-state; 99.7th pctl joint multi-state ($10{,}000$ draws)
  & Helps & $K^*$ is near-best without optimization \\
SOTA methodology (Sec.~\ref{sec:sota-methodology})
  & 7 non-adaptive strategies + 2 adaptive controls at matched budget; $K^*$ first overall ($p < 10^{-4}$)
  & Helps & Quantitative superiority over all tested alternatives \\
Hybrid $K^*$+adaptive (Sec.~\ref{sec:hybrid-test})
  & Adding adaptive rounds to $K^*$ gives $\pm 0.001$; no marginal benefit
  & Neutral & $K^*$ already near-optimal as single batch \\
Reconstructor dependence (Sec.~\ref{sec:significance})
  & Under linear inversion, random outperforms $K^*$ by 27\%
  & Limitation & Advantage requires nonlinear reconstruction \\
Shot allocation (Sec.~\ref{sec:shot-alloc})
  & Non-uniform shot weighting gives $< 0.01$ improvement
  & Neutral & Uniform 1{,}000 shots/basis is sufficient \\
Scaling to $n = 5$ (Sec.~\ref{sec:scaling-n5})
  & Allocation fraction $\to 1$; intra-weight residual $\Delta = 0.004$
  & Helps & Weight allocation nearly sufficient beyond sweet spot \\
MLE convergence (Sec.~\ref{sec:convergence})
  & Converges in $< 200$ iterations; no sensitivity to initialization
  & Helps & Reconstruction is reproducible \\
GHZ hedge insensitivity (Sec.~\ref{sec:hedge-scan})
  & $F \approx 0.50$ for all $\zeta \in [0, 0.2]$
  & Limitation & GHZ deficit is structural, not regularization \\
Extended limitations (Sec.~\ref{sec:extended-limitations})
  & Calibration drift, shot sensitivity, Bayesian alternatives documented
  & Neutral & Full robustness evidence \\
Longevity analysis (Sec.~\ref{sm:longevity})
  & Crossover at $\varepsilon_{\mathrm{eff}} \approx 10^{-3}$; 3 structural advantages persist post-FT
  & Helps & Maps $K^*$ operating envelope through fault tolerance \\
\end{tabular}
\end{ruledtabular}
\end{table*}

The supplemental analyses establish that the $K^*$ framework:
\begin{enumerate}
\item Achieves a $4.7\times$ circuit compression (29 circuits, not
137), recovering 76.9\% of operators, the single most impactful
practical finding (main text, Appendix~D).
\item Provides \emph{informationally complete} 1- and 2-body
observable recovery ($F = 0.985$ and $0.771$ respectively),
reframing $K^*$ as an RDM-quality method rather than a full-state
tomography method (main text, Appendix~E).
\item Operates in a regime where standard measurement-design
criteria (Fisher information, frame potential, coherence) are
uninformative; the advantage is a nonlinear reconstruction
phenomenon mediated by positivity-constrained MLE
(main text, Appendix~F).
\item Extends to qudits with a modest threshold relaxation
($K = q^2 + 1$) restoring low-weight saturation
(Sec.~\ref{sec:qudit}).
\item Outperforms all tested alternatives on hardware:
$\Delta F = +0.33 \pm 0.07$ (IBM, $n = 4$), $+0.25$ (Rigetti,
$n = 4$), $+0.53 \pm 0.06$ (Rigetti, $n = 8$), and ranks first
among seven non-adaptive strategies at matched budget
(Sec.~\ref{sec:optimality-test}; two adaptive baselines serve as
out-of-regime controls).
\item Derives its MLE advantage primarily from weight-class
allocation ($76 \pm 12$\% on hardware), with weight-2 completeness
identified as the dominant contribution
(Appendix~F of the main text).
\end{enumerate}

Equally important is what $K^*$ is \emph{not}:
\begin{itemize}
\item Not Fisher-optimal (all Pauli subsets of the same size are
Fisher-equivalent; main text, Appendix~F), but empirically
near-optimal under MLE (Sec.~\ref{sec:optimality-test}).
\item Not frame-optimal (all Pauli subsets have identical frame
metrics; main text, Appendix~F).
\item Not a full-state tomography method (53.3\% coverage leaves
the state underdetermined; the value proposition is local
observable recovery).
\item Not a qudit panacea at large $n$ (coverage drops to 6.5\%
for $q = 3$, $n = 4$ at strict $K^* = q^2$; however, dimensional
rigidity prescribes $2$-qutrit patches at 36.3\% coverage, and the
correct architecture is compositional, not monolithic;
simulation validates this: compositional $1$-RDM $F = 0.912$
vs.\ monolithic $0.861$ for the W~state;
Sec.~\ref{sec:qudit}).
\end{itemize}

%% ============================================================
%% MOVED APPENDICES G--L
%% ============================================================


%% ============================================================
%% Former SM S5
%% ============================================================
\section{Extension to qudit systems}
\label{sec:qudit}

\subsection{Algebraic framework}
\label{sec:qudit-framework}

Krawtchouk spectral analysis extends from qubits ($q = 2$) to qudits
of local dimension $q \geq 3$ by replacing the binary Hamming
scheme $H(n,2)$ with the $q$-ary scheme $H(n,q)$.

Generalized Pauli support classes ($q^n$ total) are indexed by
$\epsilon \in \mathbb{Z}_q^n$.
A Gram matrix measuring spectral distinguishability at cutoff $K$
is the character sum
\begin{equation}
G_{\epsilon\epsilon'}(K)
= \sum_{\substack{m \in \mathbb{Z}^n \\ |m|^2 \leq K}}
   \omega^{\bar{m} \cdot (\epsilon \ominus \epsilon')},
\label{eq:gram_q}
\end{equation}
where $\omega = e^{2\pi i/q}$, $\bar{m} = m \bmod q$, and
$\epsilon \ominus \epsilon'$ denotes coordinate-wise subtraction
in $\mathbb{Z}_q$.

\pclaim{env:prop:spectral_q}{proved}{*}{prop:spectral_q}%
\begin{proposition}[Spectral decomposition (cf.\ main text Lemma~5)]
\label{prop:spectral_q}
Let $c_w(K) = |\{m \in \mathbb{Z}^n : |m|^2 \leq K,\;
w(\bar{m}) = w\}|$.

\noindent\emph{(i) All $q$.}\;
$G(K)$ is weight-block-diagonal ($S_n$-invariant), with
mean eigenvalue (trace divided by block dimension)
\begin{equation}
\lambda_w(K) = \frac{q^n \, c_w(K)}{\binom{n}{w}(q-1)^w}
\label{eq:eigenvalue_q}
\end{equation}
on each weight-$w$ block.

\noindent\emph{(ii) $q \le 3$ only.}\;
Each weight-$w$ block is scalar ($\lambda_w$ is the unique
eigenvalue), $G(K)$ lies in the Bose--Mesner algebra of $H(n,q)$,
and decomposes as
$G(K) = \sum_{w=0}^{n} \lambda_w(K)\, E_w$,
where $E_w$ are the primitive idempotents of $H(n,q)$ with
multiplicity $\binom{n}{w}(q-1)^w$.
This fails for $q \ge 4$.
\end{proposition}

\begin{proof}
\emph{(Iff characterization.)}
Coordinate-permutation invariance ($S_n$-symmetry) of the
Euclidean norm makes $G(K)$ weight-block-diagonal for all~$q$.
For $q \le 3$, every nonzero residue $a \in \mathbb{Z}_q \setminus
\{0\}$ has minimal symmetric lift with $|\ell(a)|^2 = 1$
(for $q = 2$: $\ell(1) = \pm 1$; for $q = 3$: $\ell(1) = 1$,
$\ell(2) = -1$), so the Euclidean cutoff is additionally
Hamming-distance-invariant.  Therefore $G(K)$ lies in the
Bose--Mesner algebra of $H(n,q)$, and the $q$-ary Krawtchouk
transform diagonalises it in the idempotent basis
$\{E_w\}_{w=0}^{n}$ with eigenvalues
\[
  \lambda_w(K) \;=\; \sum_{h=0}^{K} K_h(w;\,n,q)\,G^{(h)}(K),
\]
where $K_h(w;\,n,q)$ are the $q$-ary Krawtchouk polynomials and
$G^{(h)}$ are the distance values of $G(K)$.
Self-duality of the $q$-ary Krawtchouk matrix
($\sum_h K_h(w) K_h(w') = q^n \delta_{ww'}$) yields the closed-form
spectral decomposition~\eqref{eq:eigenvalue_q}.
For $q \ge 4$, residue~$1$ lifts to $|\ell(1)|^2 = 1$ while
residue~$2$ lifts to $|\ell(2)|^2 = 4$; these distinct squared-lift
norms break Hamming-distance invariance.
At $n = 1$, $K = 1$, the off-diagonal entries of~$G$ within
Hamming distance class~$1$ are not all equal
($G_{0,1} = 1 + 2\cos(2\pi/q) \ne G_{0,2} = 1 + 2\cos(4\pi/q)$
for $q = 4$: values $1$ vs $-1$), so
$G(K) \notin \mathrm{BM}(H(1,q))$.
\end{proof}

\begin{remark}[Saturation thresholds]
For $q = 2$, $K = 5$ saturates weight classes $1$ and $2$ at every
$n \ge 4$ (weight-$1$ directly: $c_1(5) = 2n(2n{-}1) > 3n = A_1$;
weight-$2$ after redistribution of excess weight-$0,1$ budget:
e.g.\ at $n = 6$, $c_1(5) = 132 > 18 = A_1$ and
$c_2(5) = 60 < 135 = A_2$, but the $126$-unit excess from
weights~$0,1$ brings the redistributed budget to $186 > 135$;
see Sec.~\ref{sec:dim-proof} for the full proof).
Full saturation to $\lfloor n/2 \rfloor$ holds only for
$n \le 5$, where $\lfloor n/2 \rfloor \le 2$ and the weight-$1,2$
guarantee suffices.
For $n \ge 6$, weight-$3$ saturation requires
$c_3(K) \ge A_3 = \binom{n}{3} 27$, which is not achieved at $K = 5$
(e.g.\ at $n = 6$, $c_3(5) = 160 < 540 = A_3$); the monolithic
full-saturation cutoff $K_{\mathrm{full}}(n)$ grows with~$n$, and
the compositional architecture replaces the monolithic budget
(Section on compositionality).
For $q \ge 3$, $K^* = K_{\mathrm{mass}} = q^2$ is the spectral
mass threshold (Proposition~\ref{prop:mass-threshold}), where
the mod-$q$ invisible vectors enter and the eigenvalue
distribution undergoes a structural transition; weight
saturation is partial at this cutoff and full saturation
requires $K_{\mathrm{sat}} > q^2$.
These thresholds are verified by explicit enumeration of
$\lambda_w(K)$ and $M_w(K)$ at $n = 2,\ldots,9$
(Table~\ref{tab:q2-progression}).
\end{remark}

\pclaim{env:cor:support}{proved}{q>=2}{cor:support}%
\begin{corollary}[Support-completeness (cf.\ main text Lemma~6, $q = 2$)]
\label{cor:support}
If $\lambda_w(K) > 0$ for all $w = 0, \ldots, n$, then every
support class of the $n$-qudit generalized Pauli algebra is
populated and $\operatorname{rank} G(K) = q^n$.
\end{corollary}

\begin{proof}
$\lambda_w > 0$ implies $c_w > 0$ (by the eigenvalue
formula $\lambda_w = q^n c_w / \binom{n}{w}(q{-}1)^w$),
and coordinate-permutation symmetry ($S_n$-invariance)
extends to all $\binom{n}{w}$ patterns of weight~$w$.
Hence every support class is populated and $\operatorname{rank} G(K) = q^n$.
\end{proof}

\subsection{Eigenvalue tables}
\label{sec:eigenvalues}

\paragraph{Qubit case: $q = 2$, $n = 4$, $K^* = 5$.}
Table~\ref{tab:q2-progression} shows the shell-by-shell
accumulation of the parity-weight counts $c_w(K)$ and
the redistributed budget $M_w(K)$ from $K = 1$ to $K^* = 5$.
At $K \le 3$, $c_{K+1}(K) = 0$ (bold): the primitive idempotent
$E_{K+1}$ provides a Delsarte dual certificate of spectral
incompleteness.
At $K = 4$, all $c_w > 0$ (full Gram rank), but greedy
redistribution yields $M_2 = 28 < 54 = A_2$: weight-$2$
is not saturated and Theorem~2(i) gives worst-case $F \le 1/2$.
At $K^* = 5$, the $48$ weight-$1$ vectors from the $|m|^2 = 5$ shell
push the redistributed budget to $M_w = \{1, 12, 54, 54, 16\}$,
achieving weight saturation without manual reallocation.

\begin{table}[ht!]
\caption{Spectral progression at $q = 2$, $n = 4$.
$c_w(K)$: cumulative parity-weight counts.
$M_w$: redistributed budget (greedy, lowest weight first).
$A_w = \binom{4}{w}3^w$: available operators per weight.
Bold zeros indicate $c_w = 0$ (Delsarte obstruction).}
\label{tab:q2-progression}
\begin{ruledtabular}
\begin{tabular}{crrrrrrl}
$K$ & $c_0$ & $c_1$ & $c_2$ & $c_3$ & $c_4$ & $N_4(K)$ & Status \\
\hline
1 & 1 & 8 & $\mathbf{0}$ & $\mathbf{0}$ & $\mathbf{0}$ & 9 & $E_2 \in \ker G$ \\
2 & 1 & 8 & 24 & $\mathbf{0}$ & $\mathbf{0}$ & 33 & $E_3 \in \ker G$ \\
3 & 1 & 8 & 24 & 32 & $\mathbf{0}$ & 65 & $E_4 \in \ker G$ \\
4 & 9 & 8 & 24 & 32 & 16 & 89 & $\mu_2 < 1$ \\
\textbf{5} & \textbf{9} & \textbf{56} & \textbf{24} & \textbf{32} & \textbf{16} & \textbf{137} & \textbf{saturated} \\
\hline
$A_w$ & 1 & 12 & 54 & 108 & 81 & 256 & \\
$M_w(5)$ & 1 & 12 & 54 & 54 & 16 & 137 & \\
\end{tabular}
\end{ruledtabular}
\end{table}

\paragraph{Qutrit case: $q = 3$, $n = 4$, $K^* = 9$.}
Exhaustive enumeration yields $N_4(9) = 425$ lattice vectors.

\begin{table}[ht!]
\caption{Eigenvalue decomposition at $q = 3$, $n = 4$, $K^* = 9$.}
\label{tab:q3}
\begin{ruledtabular}
\begin{tabular}{crrrrr}
$w$ & $\binom{4}{w}2^w$ & $c_w$ & $\lambda_w$ &
$\lambda_w \binom{4}{w}2^w$ & \% tr$(G)$ \\
\hline
0 & 1  & 9   & 729  & 729    & 2.1  \\
1 & 8  & 16  & 162  & 1296   & 3.8  \\
2 & 24 & 96  & 324  & 7776   & 22.6 \\
3 & 32 & 224 & 567  & 18144  & 52.7 \\
4 & 16 & 80  & 405  & 6480   & 18.8 \\
\hline
  & 81 & 425 &      & 34425  & 100  \\
\end{tabular}
\end{ruledtabular}
\end{table}

Weight-3 dominates (52.7\%), in contrast to the qubit
case where weight-1 dominates (40.9\%).

\paragraph{Ququint case: $q = 5$, $n = 4$, $K^* = 25$.}
$N_4(25) = 3121$ lattice vectors.

\begin{table}[ht!]
\caption{Eigenvalue decomposition at $q = 5$, $n = 4$, $K^* = 25$.}
\label{tab:q5}
\begin{ruledtabular}
\begin{tabular}{crrrrr}
$w$ & $\binom{4}{w}4^w$ & $c_w$ & $\lambda_w$ &
$\lambda_w \binom{4}{w}4^w$ & \% tr$(G)$ \\
\hline
0 & 1   & 9    & 5625.0  & 5625    & 0.3  \\
1 & 16  & 32   & 1250.0  & 20000   & 1.0  \\
2 & 96  & 360  & 2343.8  & 225000  & 11.5 \\
3 & 256 & 1216 & 2968.8  & 760000  & 39.0 \\
4 & 256 & 1504 & 3671.9  & 940000  & 48.2 \\
\hline
  & 625 & 3121 &         & 1950625 & 100  \\
\end{tabular}
\end{ruledtabular}
\end{table}

\subsection{Characterization of spectral thresholds}
\label{sec:threshold}

We prove two threshold results analytically and
identify the operational threshold for qubits.

\pclaim{env:prop:spec-complete}{proved}{n>=1,q>=2}{prop:spec-complete}%
\begin{proposition}[Spectral completeness threshold]
\label{prop:spec-complete}
For any $q \ge 2$ and $n \ge 1$, the Gram matrix $G(K)$
of $H(n,q)$ achieves spectral completeness ($c_w(K) > 0$
for all $w = 0, \ldots, n$) at $K = n$.
This threshold is tight: $c_n(n{-}1) = 0$.
\end{proposition}
\begin{proof}
\emph{Sufficiency.}
For each weight $w \ge 1$, the vector
$m = (\underbrace{1,\ldots,1}_{w},\, 0,\ldots,0)$ has
$|m|^2 = w \le n$ and $w(\bar{m}) = w$ (since $1 \not\equiv 0
\pmod{q}$ for $q \ge 2$).  Hence $c_w(n) \ge 1$ for $w \ge 1$.
For $w = 0$, the origin gives $c_0(K) \ge 1$ for all $K \ge 0$.

\emph{Tightness.}
Any lattice vector of parity weight~$n$ has all $n$
coordinates nonzero mod~$q$, hence $|m|^2 \ge n$.
Therefore $c_n(n{-}1) = 0$, and the idempotent $E_n$
provides a Delsarte dual certificate at $K = n{-}1$.
\end{proof}

\pclaim{env:prop:mass-threshold}{proved}{n>=1,q>=2}{prop:mass-threshold}%
\begin{proposition}[Spectral mass threshold]
\label{prop:mass-threshold}
For any $q \ge 2$ and $n \ge 1$, the first non-origin
lattice points invisible to mod-$q$ reduction appear at
$|m|^2 = q^2$.  Specifically:
\begin{enumerate}
\item $c_0(q^2{-}1) = 1$ (origin only).
\item $c_0(q^2) = 1 + 2n$.
\item The jump is universal: it equals $2n$ for
  all $q \ge 2$.
\end{enumerate}
\end{proposition}
\begin{proof}
A nonzero vector $m \in \mathbb{Z}^n$ has parity weight
$w(\bar{m}) = 0$ iff every coordinate is divisible by~$q$.
The smallest such vector has the form $(\pm q, 0, \ldots, 0)$
with $|m|^2 = q^2$.
Below $|m|^2 = q^2$, no nonzero vector can have all
coordinates divisible by~$q$:
if $m_i = q\, a_i$ with $a \neq 0$, then
$|m|^2 = q^2 |a|^2 \ge q^2$.
At $|m|^2 = q^2$, the solutions are exactly the $2n$
vectors $\pm q\, e_i$ ($i = 1, \ldots, n$), giving
$c_0(q^2) = 1 + 2n$.
\end{proof}

\paragraph{Full rank is necessary but not sufficient.}
By Proposition~\ref{prop:spec-complete},
$\operatorname{rank} G(K) = q^n$ is first achieved at $K = n$,
independent of $q$.  However,
at $K = n$ the eigenvalue structure is flat (no useful weight-class
discrimination).  For $q = 2$, the weight-1 sector carries only
9\% of the eigenvalue mass at $K = 4$; $K^* = 5$ raises this to
40.9\%.

\paragraph{The $|m|^2 = q^2$ shell and mod-$q$ invisible vectors.}
By Proposition~\ref{prop:mass-threshold}, the vectors
$m = (\pm q, 0, \ldots, 0)$ and permutations at $|m|^2 = q^2$
are the first non-origin lattice points invisible to
mod-$q$ reduction.  Their inclusion increases
$c_0$ from 1 to $1 + 2n = 9$ at $n = 4$, universally across $q$.

\begin{table}[ht!]
\caption{Effect of the $k = q^2$ shell on $c_0$ at $n = 4$.}
\label{tab:shell}
\begin{ruledtabular}
\begin{tabular}{ccccc}
$q$ & $K_{\mathrm{mass}} = q^2$ & $c_0(K_{\mathrm{mass}}{-}1)$ & $c_0(K_{\mathrm{mass}})$ &
shell $r_4(q^2)$ \\
\hline
2 & 4  & 1 & 9   & 24  \\
3 & 9  & 1 & 9   & 104 \\
5 & 25 & 1 & 9   & 248 \\
7 & 49 & 1 & 9   & 456 \\
\end{tabular}
\end{ruledtabular}
\end{table}

\paragraph{Shell-by-shell eigenvalue mass progression.}

Table~\ref{tab:progression} shows the accumulation of $c_w$ and
eigenvalue mass for $q = 3$, $n = 4$.  Three features emerge:
(i)~full rank at $K = 4$ with flat eigenvalues;
(ii)~weight-3 dominance from $K = 6$, locking in at $K = 9$;
(iii)~$c_0$ jumping from 1 to 9 at $K = 9$ (mod-$q$ invisible
vectors).

\begin{table}[ht!]
\caption{Eigenvalue mass progression at $q = 3$, $n = 4$.}
\label{tab:progression}
\begin{ruledtabular}
\begin{tabular}{crrrrrrrrrr}
 & \multicolumn{5}{c}{$c_w(K)$}
 & \multicolumn{5}{c}{\% of tr$(G)$} \\
\cline{2-6}\cline{7-11}
$K$ & $w{=}0$ & $1$ & $2$ & $3$ & $4$
    & $w{=}0$ & $1$ & $2$ & $3$ & $4$ \\
\hline
1 & 1 &  8 &   0 &   0 &  0
  & 11.1 & 88.9 & $-$ & $-$ & $-$ \\
2 & 1 &  8 &  24 &   0 &  0
  & 3.0 & 24.2 & 72.7 & $-$ & $-$ \\
3 & 1 &  8 &  24 &  32 &  0
  & 1.5 & 12.3 & 36.9 & 49.2 & $-$ \\
4 & 1 & 16 &  24 &  32 & 16
  &  1.1 & 18.0 & 27.0 & 36.0 & 18.0 \\
5 & 1 & 16 &  72 &  32 & 16
  &  0.7 & 11.7 & 52.6 & 23.4 & 11.7 \\
6 & 1 & 16 &  72 & 128 & 16
  &  0.4 & 6.9 & 30.9 & 54.9 & 6.9 \\
7 & 1 & 16 &  72 & 128 & 80
  &  0.3 & 5.4 & 24.2 & 43.1 & 26.9 \\
8 & 1 & 16 &  96 & 128 & 80
  &  0.3 & 5.0 & 29.9 & 39.9 & 24.9 \\
9 & 9 & 16 &  96 & 224 & 80
  &  2.1 & 3.8 & 22.6 & 52.7 & 18.8 \\
\end{tabular}
\end{ruledtabular}
\end{table}

Three thresholds emerge from
Propositions~\ref{prop:spec-complete}--\ref{prop:mass-threshold}:
(i)~spectral completeness at $K = n$ (all $c_w > 0$; proved);
(ii)~the spectral mass transition at $K_{\mathrm{mass}} = q^2$
where $c_0$ jumps from~$1$ to $1 + 2n$ (proved);
(iii)~operational weight saturation ($M_w = A_w$ for
$w = 1, 2$) at $K = 5$ for $q = 2$, all $n \ge 4$ (proved);
full saturation to $\lfloor n/2 \rfloor$ holds for $n \le 5$;
for $q = 3$, full saturation requires
$K_\mathrm{sat} = 14$ (verified computationally).
For $q = 2$, thresholds (ii) and (iii) are adjacent:
$K_{\mathrm{mass}} = 4$, $K^* = 5$.
For $q \ge 3$, $K^*$ in this paper denotes
$K_{\mathrm{mass}} = q^2$ (threshold~(ii)); full weight saturation
requires $K_{\mathrm{sat}} > q^2$ and several additional shells;
the compositional architecture (Sec.~\ref{sec:qudit-rdm})
compensates by decomposing into smaller patches.
Whether $K_\mathrm{sat}$ admits a closed-form expression in
$q$ and $n$ for $q \ge 3$ is an open question.

\subsection{Weight-class saturation and coverage scaling}
\label{sec:saturation}

A key structural difference between qubits and qudits emerges
in weight-class saturation.  For $q = 2$, $K = 5$ saturates
weight classes $w \in \{1, 2\}$ for every $n \ge 4$.
At $K = 5$, each nonzero coordinate satisfies $|m_i| \le 2$
(since $|m_i| \ge 3$ would require $|m|^2 \ge 9 > 5$),
and $2 \equiv 0 \pmod{2}$, so $\pm 2$ entries do not
increase parity weight.
For parity weight~$w = 1$: one odd coordinate ($\pm 1$)
plus up to one even companion ($\pm 2$, contributing
$|m_i|^2 = 4$) gives
$c_1(5) = 2n + 4n(n{-}1) = 2n(2n{-}1)$, while
$A_1 = 3n$.
Hence $c_1 / A_1 = 2(2n{-}1)/3 > 1$ for all $n \ge 1$.
For parity weight~$w = 2$: two odd coordinates ($\pm 1, \pm 1$,
$|m|^2 = 2$, leaving no budget for $\pm 2$ companions) gives
$c_2(5) = 4\binom{n}{2}$, while $A_2 = 9\binom{n}{2}$.
Here $c_2 < A_2$ (ratio $4/9$), so weight-$2$ saturation
requires redistribution: excess budget from weight-$0$
($c_0 - 1 = 2n$) and weight-$1$ ($c_1 - A_1 = n(4n{-}5)$)
flows to weight-$2$, and the combined total
$6n^2 - 5n \ge 9\binom{n}{2}$ holds for all $n \ge 1$
(Sec.~\ref{sec:dim-proof}).
Full saturation to $\lfloor n/2 \rfloor$ holds only for
$n \le 5$, where $\lfloor n/2 \rfloor \le 2$.
For $n \ge 6$, weight-$3$ saturation is not achieved at $K = 5$:
e.g., at $n = 6$, $c_3(5) = 160$ vs $A_3 = 540$ (30\% coverage),
because lattice points with $3$ odd coordinates contribute
$|m|^2 \ge 3$ already and additional non-zero even coordinates
(needed to span all $3^3 = 27$ Pauli signs per index choice)
push $|m|^2$ above $5$.
The monolithic full-saturation cutoff $K_{\mathrm{full}}(n)$
grows with $n$ ($K_{\mathrm{full}} = 6$ at $n = 6$, $7$ at
$n = 8$), and the corresponding operator budget scales rapidly
($1{,}341$ at $n = 6$; $12{,}369$ at $n = 8$).
The compositional architecture circumvents this scaling by
keeping each patch at $n = 4$, where $K^* = 5$ achieves full
saturation with $137$~operators per patch.
For $q = 3$, partial saturation at $K = q^2$ is more severe:

\begin{table}[ht!]
\caption{Weight-class saturation at $K^* = q^2{+}1$ ($q = 2$) or
$q^2$ ($q \ge 3$) and minimum $K$ for full saturation ($n = 4$).}
\label{tab:saturation}
\begin{ruledtabular}
\begin{tabular}{ccccccc}
$q$ & Wt & Total & $K^*$ & At $K^*$ & Cov.\ & $K_\mathrm{sat}$ \\
\hline
2 & 1 &  12 &  5 &  12 & 100\% & 5 \\
2 & 2 &  54 &  5 &  54 & 100\% & 5 \\
\hline
3 & 1 &  32 &  9 &  16 &  50\% & 10 \\
3 & 2 & 384 &  9 &  96 &  25\% & 14 \\
\end{tabular}
\end{ruledtabular}
\end{table}

The paper's algorithm redistributes excess budget from
over-saturated classes ($c_0 = 9$ but only 1 identity operator
exists $\to$ 8 units redistribute to weight-1), raising effective
weight-1 coverage to $\sim$75\% (24/32), still short of
saturation.

Weight-1 saturation requires only $K = 10$ ($= q^2 + 1$ for $q = 3$).
Weight-2 saturation requires $K_{\mathrm{sat}} = 14 < 2q^2 = 18$, well within
the spectral framework.  For practical qudit tomography, scanning
$K \in [q^2, K_{\mathrm{sat}}]$ (here $[9, 14]$) restores the
low-weight saturation that drives the qubit advantage.

The overall coverage ratio scales unfavorably, as shown by the
lattice point count $N_4(K)$ for $q = 3$, $n = 4$:

\begin{table}[ht!]
\caption{Lattice point count and coverage for $q = 3$, $n = 4$ as
a function of cutoff $K$.}
\label{tab:coverage-scan}
\begin{ruledtabular}
\begin{tabular}{crr}
$K$ & $N_4(K)$ & Coverage (\%) \\
\hline
9 ($= q^2$) & 425 & 6.5 \\
15 & 1{,}257 & 19.2 \\
20 & 2{,}041 & 31.1 \\
25 & 3{,}121 & 47.6 \\
\textbf{26} & \textbf{3{,}457} & \textbf{52.7} \\
30 & 4{,}785 & 72.9 \\
\end{tabular}
\end{ruledtabular}
\end{table}

Reaching 50\% coverage requires $K \approx 26$, compared to
$K^* = 5$ for the qubit case.  At $K = 26$, the eigenvalue
structure is:

\begin{table}[ht!]
\caption{Eigenvalue decomposition at $q = 3$, $n = 4$, $K = 26$
(approximately 50\% coverage).}
\label{tab:q3-k26}
\begin{ruledtabular}
\begin{tabular}{crrrr}
$w$ & $c_w$ & Mult.\ & $\lambda_w$ & Mass (\%) \\
\hline
0 & 33  & 1  & 2{,}673 & 1.0 \\
1 & 368 & 8  & 3{,}726 & 10.6 \\
2 & 1{,}200 & 24 & 4{,}050 & 34.7 \\
3 & 1{,}152 & 32 & 2{,}916 & 33.3 \\
4 & 704 & 16 & 3{,}564 & 20.4 \\
\end{tabular}
\end{ruledtabular}
\end{table}

At comparable coverage, the qudit eigenvalue mass is more evenly
distributed across weight classes than the qubit case (where
weight-1 dominates at 40.9\%), reflecting the larger operator space
per weight class.  The extension is most immediately applicable to
small qudit registers ($n = 2$--$3$), matching current hardware
scale.

The 6.5\% coverage at $q = 3$, $n = 4$ need not represent a
fundamental limitation.  The Krawtchouk allocation condition ($d^* = 2$ for
$q \geq 3$; main text Sec.~II\,A) suggests that qutrit registers
should be decomposed into overlapping $2$-qutrit patches, where
coverage is $36.3$\%, mirroring the qubit compositional
architecture ($4$-qubit patches at $d^* = 4$).  This
compositional qutrit prescription is validated in simulation:
overlapping $2$-qutrit patches achieve $1$-RDM fidelity
$F = 0.912$ for the W~state vs.\ $0.861$ monolithic
(Sec.~\ref{sec:qudit-sim}).

\subsection{Simulation evidence}
\label{sec:qudit-sim}

The $q^{2n}$ Heisenberg--Weyl operators
$D(\mathbf{a}, \mathbf{b}) =
X^{a_1} Z^{b_1} \otimes \cdots \otimes X^{a_n} Z^{b_n}$
span the operator space for $n$ qudits of dimension~$q$, but are
\emph{not} Hermitian for $q \geq 3$.  Because tomographic
expectation values require Hermitian observables, we work in the
basis of $n$-qudit generalized Gell-Mann operators (tensor products
of single-qudit traceless Hermitian generators and identity),
yielding $q^{2n} - 1$ nontrivial basis elements classified by
support weight $w \in \{1, \ldots, n\}$.

The structured measurement set of budget $M = N_n(K^*)$ is
constructed by saturating low-weight sectors first and filling the
remaining budget from higher-weight sectors by stride sampling,
following the qubit prescription (Sec.~II\,A of the main text).  This is
a weight-prioritized allocation; the specific intra-weight ordering
has not been optimized for the qudit case.

We compare three arms: (i)~$K^*$-structured ($M$ operators selected
by weight-prioritized Krawtchouk allocation), (ii)~allocation-random
($M$ operators with the same per-weight budget but random
intra-weight selection), and (iii)~uniform-random ($M$ operators
drawn uniformly from the $q^{2n} - 1$ nontrivial elements).  All
arms use robust MLE reconstruction (matching the qubit pipeline).
Noise model: uniform depolarizing noise at $p = 0.03$ with
$N_s = 1000$ shots per operator, 10 independent seeds.

\begin{table*}
\caption{Three-arm qutrit simulation at $(q, n) = (3, 2)$,
$M = 29$ operators (Gell-Mann basis), depolarizing noise
$p = 0.03$, 1000 shots, robust MLE, 10 seeds.}
\label{tab:sim}
\begin{ruledtabular}
\begin{tabular}{lccccc}
State & $F_{K^*}$ & $F_\mathrm{AR}$ & $F_\mathrm{UR}$ &
$\Delta F(K^*{-}\mathrm{AR})$ & $\Delta F(K^*{-}\mathrm{UR})$ \\
\hline
W       & $0.504 \pm 0.017$ & $0.232 \pm 0.106$ &
$0.273 \pm 0.255$ & $+0.273$ & $+0.231$ \\
Product & $0.954 \pm 0.008$ & $0.910 \pm 0.108$ &
$0.382 \pm 0.300$ & $+0.044$ & $+0.572$ \\
GHZ     & $0.253 \pm 0.025$ & $0.272 \pm 0.225$ &
$0.171 \pm 0.098$ & $-0.020$ & $+0.082$ \\
\end{tabular}
\end{ruledtabular}
\end{table*}

Three qualitative features of the qubit results carry over:

\paragraph{Stability.}
The $K^*$ arm exhibits low variance ($\sigma \leq 0.025$),
while allocation-random is $6\times$ more volatile ($\sigma$ up
to $0.225$) and uniform-random exceeds $\sigma = 0.25$.

\paragraph{Structured advantage for locally informative states.}
W and product states show positive $\Delta F(K^*{-}\mathrm{AR})$.
The W-state advantage ($+0.273$) is the largest, matching the
qubit pattern where states with nontrivial single-site marginals
benefit most from Krawtchouk weight allocation.

\paragraph{GHZ pattern.}
The qutrit GHZ state shows
$\Delta F(K^*{-}\mathrm{AR}) = -0.020$, not statistically
significant, but directionally consistent with the eigenvalue-mass
mismatch: GHZ information resides in high-weight global correlators
that are subsampled by the structured allocation.

\subsection{Absolute fidelity and underdetermination}

At $(q,n) = (3,2)$, full-state fidelities remain moderate
($F_{K^*} = 0.504$ for~W) because 29 operators cannot fully
determine a $9$-dimensional state.  However, the operationally
relevant metric is local-observable accuracy, not full-state
reconstruction (see Sec.~\ref{sec:qudit-rdm} below).

\subsection{Fidelity versus budget}
\label{sec:budget}

A budget sweep from $M = 5$ to 80 for $(q,n) = (3,2)$ under
linear inversion confirmed that the structured advantage persists
across budgets and vanishes at full tomography, ruling out
systematic bias.  The qualitative pattern matches the qubit case
(Appendix~F of the main text): product states show a monotonic
structured advantage that vanishes at full tomography, while the
W~state shows an initial deficit at very low $M$ (weight-$1$
saturation wastes budget before weight-$2$ fills) followed by
structured recovery as $M$ increases.  The MLE reconstructor is
expected to sharpen these trends, as it does for qubits
(the ablation table (main text, Appendix~F)).

\subsection{RDM implications for qudits}
\label{sec:qudit-rdm}

The qubit RDM results (Appendix~E of the main text) have direct
implications for the qudit extension.  At $K = q^2 + 1$ (weight-1
saturation), all single-qudit RDMs are fully determined.  At
$K = 14$ (weight-2 saturation for $q = 3$, $n = 4$), all
two-qudit RDMs are complete.  The most immediate experimental
target is two qutrits ($q = 3$, $n = 2$): 29 operators out of 81,
matching current qutrit hardware scale.

Simulation confirms this: at $(q,n) = (3,2)$ with 29 $K^*$
operators, the W-state $1$-RDM fidelity is
$F_\mathrm{1\text{-}RDM} = 0.972 \pm 0.028$, even though full-state
fidelity is only $0.504$.  The GHZ state achieves
$F_\mathrm{1\text{-}RDM} = 0.915 \pm 0.044$.  Per-qutrit properties
are recovered at high accuracy despite the underdetermined global
reconstruction, validating the Krawtchouk allocation for
certification tasks where local observables are the target.

The compositional architecture extends this to larger registers:
overlapping $2$-qutrit patches on a $4$-qutrit system achieve
$1$-RDM fidelity $F = 0.912$ (W~state) vs.\ $0.861$ for
monolithic $4$-qutrit $K^*$, confirming the $d^* = 2$ patch
prescription.


%% ============================================================
%% Former SM S6
%% ============================================================
\section{Dimensional dependence: proof details}
\label{sec:dim-proof}

This section provides the full low-weight saturation proof and
reconstructor cross-checks supporting the dimensional progression
$\Delta(S{-}AR): 0.573 \to 0.331 \to 0.100 \to 0.042$ reported in
Sec.~V\,C of the main text.

\subsection{Low-weight saturation theorem}

A lattice vector $v \in \mathbb{Z}^n$ with $\lVert v\rVert^2 \leq K$
has \emph{parity weight} $w = |\{i : v_i \not\equiv 0 \pmod{2}\}|$;
the parity-weight budget $c_w$ counts such vectors.
For parity-weight~$1$, exactly one coordinate is odd
($v_i = \pm 1$, since $|\pm 3|^2 = 9 > 5$) and all others
are even ($0$ or $\pm 2$).  The remaining norm budget is
$K - 1$, and each even coordinate contributes $0$ or $4$;
thus each non-odd position can be $0$ or $\pm 2$ (if
$K - 1 \geq 4$, i.e.\ $K \geq 5$), giving at most one
$\pm 2$ companion.  This yields:
$c_1 = 2n + 4n(n{-}1) = 2n(2n{-}1)$ at $K = 5$
(the $2n$ pure single-odd vectors plus $4n(n{-}1)$ mixed
$(\pm 1, \pm 2)$ vectors; note $(\pm 2, 0, \ldots)$ has
parity weight~$0$, not~$1$).
At $n = 4$: $c_1 = 56$, matching the main-text enumeration.
Since $A_1 = 3n$, the inequality $c_1 > A_1$ holds for all
$n \geq 2$.
Similarly, at $K^* \geq 5$ the weight-$2$ budget satisfies
$c_2 \geq 4\binom{n}{2}$
(the pairs $(\pm 1, \pm 1)$ contribute $4$ vectors per
coordinate pair; note that $(\pm 1, \pm 2)$ has
parity weight~$1$ since $2 \equiv 0 \pmod{2}$).
Since the $\binom{n}{2}$ weight-$2$ parity patterns are
each realized by $(\pm 1, \pm 1)$, all weight-$2$ support
classes are covered.
Weight-$1$ saturation follows from
$c_1 = 2n(2n{-}1) > 3n = A_1$ for all $n \geq 2$;
weight-$2$ saturation follows after redistribution:
the excess from weight-$0$ and weight-$1$
($c_0 + c_1 - 1 - A_1 = 4n^2 - 5n + 8$ at $n = 4$;
$\geq 2n(2n{-}1) - 3n = n(4n{-}5)$ in general)
combined with $c_2 = 4\binom{n}{2}$ exceeds
$A_2 = 9\binom{n}{2}$ for all $n \geq 4$.
Therefore, for $n \geq 4$ the $K^*$ redistribution includes
\emph{every} weight-$0$, weight-$1$, and weight-$2$ Pauli operator:
the structured and weight-allocated random sets are identical at these
weights.  Since weights $0$--$2$ carry all single- and two-body
correlation information, the residual $\Delta(S{-}AR)$ arises entirely
from weight $\geq 3$ differences.
For $n = 2, 3$ (where the spectral-completeness cutoff
$K = n$ is below the formal weight-saturation threshold
$K^* = 5$), weight-$1$ is not
fully covered, so selection freedom exists at all
weights including the informationally dominant weight-$1$ sector,
explaining the large $\Delta(S{-}AR)$.
The mechanism is the transition from the overcomplete
regime ($K = n$) to the weight-saturation threshold
$K^* = 5$ at $n = 4$; the Krawtchouk allocation condition
($2n = 2^{n-1}$, satisfied uniquely at $n = 4$) determines
\emph{why} the transition occurs at this dimension, via the
weight-class decomposition
$N(K^*) = 1 + 12 + 54 + 54 + 16 = 137$.

\subsection{Reconstructor robustness check}

The $n = 2$ data point uses the $R\rho R$ MLE that exhibits
identity-weight dominance at $n = 2$ (main text, Sec.~V\,A).
To verify the progression is not an MLE artifact, we reran the
three-arm simulation with linear inversion (Pauli expansion
$\rho = d^{-1}\sum_i b_i P_i$ followed by projection to the nearest
physical state) and projected least squares (PLS).

Under linear inversion,
$\Delta(S{-}AR) = +0.128$ at $n = 2$, $+0.003$ at $n = 3$, and
$+0.010$ at $n = 4$; under PLS,
$+0.127$ at $n = 2$, $+0.003$ at $n = 3$, $+0.008$ at $n = 4$, and
$+0.005$ at $n = 5$.
In all three reconstructors the $n = 2$ value remains the largest,
consistent with intra-weight selection mattering most in the
high-coverage regime.  However, both non-MLE progressions are
non-monotonic ($n{=}3 < n{=}4$), so the smooth decrease
$n = 2 > n = 3 > n = 4 > n = 5$ seen under MLE is
reconstructor-specific; the monotonic decrease through the $n = 4$
sweet spot (where the Krawtchouk allocation condition holds) is
supported by MLE but not by linear inversion or PLS.  Absolute fidelities under linear inversion and PLS are lower
($F \sim 0.23$--$0.60$) because they lack the positivity enforcement
of iterative MLE, making them poor reconstructors for the
underdetermined regime but valid cross-checks of the dimensional trend.


%% ============================================================
%% Former SM S7
%% ============================================================
\section{Intra-weight ordering robustness}
\label{sec:ordering}

The lexicographic stride used throughout this paper was fixed before
any hardware experiment and never modified.  The following analysis
tests six alternative orderings \emph{post hoc} in simulation to
quantify sensitivity; no ordering was selected based on fidelity
performance.

The intra-weight operator selection at weights~$3$ and~$4$ uses a
lexicographic stride over Pauli strings (Sec.~II\,A of the main text),
which produces an asymmetric support-pattern distribution: e.g.,
support pattern $(0,1,1,1)$ receives all $27$ operators while
$(1,1,1,0)$ receives $7$.  We tested seven orderings with identical
weight budgets $\{1, 12, 54, 54, 16\}$ (W~state, depolarizing
$p = 0.03$, readout $2$\%, $1{,}000$~shots, $10$~noise draws):

\begin{table}[ht!]
\caption{Intra-weight ordering sensitivity.  All orderings share
the weight budget $\{1, 12, 54, 54, 16\}$; only the selection of
\emph{which} operators within each weight class differs.  ``Overlap''
is the fraction of operators shared with the default.}
\label{tab:ordering}
\begin{ruledtabular}
\begin{tabular}{lccl}
Ordering & $F$ (mean$\pm$std) & Overlap & Cluster \\
\hline
random(7777)              & $0.817 \pm 0.030$ & 78.8\% & typical \\
lexicographic (default)   & $0.809 \pm 0.072$ & 100\%  & typical \\
support\_first            & $0.805 \pm 0.081$ & 98.5\% & typical \\
random(8888)              & $0.804 \pm 0.032$ & 78.1\% & typical \\
\hline
random(9999)              & $0.426 \pm 0.101$ & 78.8\% & adversarial \\
reverse                   & $0.408 \pm 0.170$ & 81.0\% & adversarial \\
antistride                & $0.310 \pm 0.172$ & 81.0\% & adversarial \\
\end{tabular}
\end{ruledtabular}
\end{table}

Results are \textbf{bimodal}: four orderings cluster at
$F \approx 0.80$--$0.82$ (median $0.807$); three cluster at
$F \approx 0.31$--$0.43$.  Our default lexicographic stride sits at
the median of the typical cluster, not at an optimized extreme.
Two of three random shuffles also land in the typical cluster,
confirming that the default is representative of generic orderings.

All three adversarial orderings share a structural property: they
concentrate operators in a few support patterns while leaving others
empty, creating reconstruction degeneracies that the MLE cannot
resolve (the same mechanism as the random arm's oscillatory
convergence; Appendix~F of the main text).  Reverse and
antistride orderings are systematic anti-correlates of the
lexicographic stride; they are adversarial by construction,
not by chance.

Note that ``parameter-free'' in the main text refers to the
weight-level budget $\{1, 12, 54, 54, 16\}$, which is identical
across all seven orderings and determined by spectral geometry.
Intra-weight ordering is a reproducibility convention; any
ordering from the typical cluster produces comparable results.

Explicitly balancing the support-pattern distribution within each
weight class ($14/14/13/13$ instead of $27/13/7/7$) does not
systematically improve fidelity ($F \in [0.76, 0.85]$ across three
seeds vs.\ $0.81$ for the default), so the residual is not primarily
explained by support-pattern imbalance.


%% ============================================================
%% Former SM S8
%% ============================================================
\section{Near-optimality among weight-respecting subsets}
\label{sec:optimality-test}

Hardware evidence establishes $K^*$'s advantage directly:
$\Delta F = +0.33 \pm 0.07$ over random Pauli sampling on W-states
(IBM \texttt{ibm\_fez}, four independent runs), corroborated at
$\Delta F = +0.25$ on Rigetti Ankaa-3, and growing to
$\Delta F = +0.53 \pm 0.06$ at $n = 8$ under the compositional
architecture.  The matched-budget comparison (seven non-adaptive
strategies plus two adaptive controls)
(Table~\ref{tab:sota-full}) shows $K^*$ ranking first across
$23$~test states in simulation.
The three-arm hardware experiment (main text, Sec.~V\,B)
attributes $76 \pm 12$\% of the advantage to the Krawtchouk weight
allocation alone.

\paragraph{Noise sensitivity of the allocation fraction.}
The ${\sim}81\%$ allocation fraction (at depol~$= 0.03$,
readout~$= 0.02$) is not an artifact of a single noise point.
A sweep over eight noise configurations
(depol~$\in \{0, 0.01, 0.02, 0.03, 0.05, 0.08, 0.10\}$,
readout~$\in \{0, 0.01, 0.02, 0.03, 0.05\}$; $20$~seeds each)
shows the allocation fraction $\Delta(AR{-}UR)/\Delta(S{-}UR)$
ranges from $54\%$ (noiseless) to $88\%$ (high noise), with
the increase reflecting that noise erodes the intra-weight
ordering advantage faster than the weight-allocation advantage.
At the manuscript's reference point (depol~$= 0.03$,
readout~$= 0.02$), the sweep gives $80.2\%$, consistent with
the reported~${\sim}81\%$.
Data: \texttt{data/auxiliary/allocation\_noise\_sensitivity.json}.

\subsection{Saturate-and-fill baseline}
\label{sec:sat-fill}

A natural heuristic is to saturate the smallest weight classes
(weights~$0$--$2$: $1 + 12 + 54 = 67$ operators) and distribute
the remaining $70$~slots proportionally by pool size among
weights~$3$ and~$4$ ($108$ and $81$ available, yielding a $40{:}30$
split).  We compare three strategies, all sharing weights~$0$--$2$
and using $M = 137$ total operators (W~state, $p = 0.03$
depolarizing, $2$\% readout, $10^3$~shots, $20$~trials, robust MLE):

\begin{center}
\begin{tabular}{lccc}
\hline
 & K$^*$ & SFK & SFP \\
 & {\small (54{:}16, det.)} & {\small (54{:}16, rand.)} & {\small (40{:}30, rand.)} \\
\hline
W & $0.817$ & $0.717$ & $0.607$ \\
GHZ & $0.806$ & $0.593$ & $0.440$ \\
Product & $0.975$ & $0.965$ & $0.958$ \\
Haar (mean) & $0.07$ & $0.07$ & $0.07$ \\
\hline
\end{tabular}
\end{center}

\noindent
SFK (``saturate-fill-$K^*$-ratio'') uses K$^*$'s $54{:}16$ split
but random intra-weight selection.  SFP (``saturate-fill-proportional'')
uses the pool-proportional $40{:}30$ split with random selection.
The K$^*$ advantage over proportional filling ($\Delta F \approx +0.21$
for W~states) decomposes as: ${\sim}\,52$\% from the weight-$3/4$
ratio (SFK vs.\ SFP: $+0.11$) and ${\sim}\,48$\% from the specific
intra-weight operator selection (K$^*$ vs.\ SFK: $+0.10$).
The Krawtchouk-prescribed allocation is not a
na\"{\i}ve saturation heuristic.


%% ============================================================
%% Former SM S10
%% ============================================================
\section{Measurement strategy comparison}
\label{sec:sota-methodology}

Table~V of the main text compares $K^*$ against eight alternative
measurement selection strategies at matched budget ($M = 137$
operators, $1{,}000$~shots each, $n = 4$ qubits).  All strategies
use the same robust MLE reconstructor and noise model ($3$\%
depolarizing, $2$\% readout per qubit), run for $20$~independent
trials on $23$~test states (product~$|{+}{+}{+}{+}\rangle$,
W, GHZ, $10$~Haar-random pure, $10$~random mixed of rank~$4$).
Source code is provided in \texttt{sota\_comparison.py} and
\texttt{sota\_adaptive.py}.

\subsection{Non-adaptive strategies}

\paragraph{Greedy D-optimal.}
Select operators one at a time to maximize $\log\det(A^\dag A + \epsilon I)$
where $A$ is the measurement matrix relating selected operators to the
full $4^n$-dimensional Pauli basis and $\epsilon = 10^{-10}$ regularizes
the initially rank-deficient Gram matrix.  At each step, the candidate
maximizing the log-determinant of the updated Gram matrix is appended.
This is the classical optimal experimental design criterion for minimizing
the volume of the confidence ellipsoid~\cite{Pukelsheim2006}.

\paragraph{Greedy E-optimal.}
Select operators to maximize the minimum eigenvalue of $A^\dag A$,
equivalent to minimizing the worst-case estimation variance.  Each
candidate is evaluated by computing all eigenvalues of the updated
Gram matrix via \texttt{numpy.linalg.eigvalsh}.

\paragraph{Greedy A-optimal.}
Select operators to minimize $\operatorname{Tr}((A^\dag A)^{-1})$
using Tikhonov regularization ($\epsilon = 10^{-8}$) when the Gram
matrix is not yet full rank.  This minimizes the average estimation
variance across all parameters.

\paragraph{Derandomized shadows.}
Greedily select tensor-product Pauli bases (elements of
$\{X, Y, Z\}^{\otimes n}$, $3^n = 81$ candidates) to cover the
maximum number of uncovered target observables.  A basis $b$
``covers'' a Pauli operator $P$ if every non-identity position in
$P$ matches the corresponding position in $b$.  The generic target
set is all $4^n - 1 = 255$ non-identity Paulis.  After all targets
are covered, remaining budget slots are filled with random bases.
Each selected basis contributes all Pauli operators it covers;
the union of measurable operators across all $137$~bases yields the
effective measurement set for reconstruction.  This simplified
protocol captures the greedy set-cover structure of Ref.~\cite{Huang2021}
without the per-observable weighting that prioritizes
high-value targets.  The omission disadvantages the shadow protocol
relative to a fully weighted implementation; the comparison therefore
provides a lower bound on derandomized-shadow performance rather
than a definitive ranking.

\paragraph{Weight-allocated random (WA-random).}
Draw $M = 137$ operators respecting the Krawtchouk weight-class
allocation ($\{1, 12, 54, 54, 16\}$ for weights $0$--$4$) but
selecting uniformly at random within each weight class.  This
isolates the contribution of weight-class structure from the
specific $K^*$ operator ordering.

\paragraph{Uniform random.}
Draw $M = 137$ operators uniformly at random from all $4^n - 1 = 255$
non-identity Paulis (fresh draw per trial).

\subsection{Adaptive strategies}

\paragraph{Adaptive 3-round.}
Split the total budget into three equal rounds ($\lfloor M/3 \rfloor$,
$\lfloor M/3 \rfloor$, remainder).  Round~1 selects operators
uniformly at random.  After each round, the accumulated expectations
are fed to the robust MLE reconstructor.  Subsequent rounds select
operators by information-gain scoring: for each candidate Pauli $P$
and current estimate $\hat\rho$, the score is
$1 - \langle P \rangle_{\hat\rho}^2$
(expected measurement variance).  Candidates are sorted by
descending score; a diversity filter discards any candidate whose
normalized Hilbert--Schmidt overlap with an already-selected operator
exceeds $0.9$.  At $n = 4$, each round has
$\lfloor 137 / 3 \rfloor = 45$ operators for $256$~parameters,
leaving every intermediate reconstruction severely underdetermined.

\paragraph{Adaptive 5-round.}
Same protocol with five rounds ($\lfloor M/5 \rfloor = 27$ per round,
remainder in round~5).  More adaptation opportunities but smaller
per-round batches ($27$~operators for $256$~parameters), making
intermediate reconstructions even less reliable.

\subsection{Results}

Table~\ref{tab:sota-full} reports per-state-class fidelities.
$K^*$ ranks first overall ($\bar{F} = 0.501$), ahead of
D-optimal ($0.479$, $\Delta = +0.022$).  Separation is largest
on the W~state ($+0.12$ over D-optimal, $+0.10$ over E-optimal)
and GHZ ($+0.28$), where $K^*$'s weight-class balance preserves
both local correlators and weight-$4$ phase correlations that greedy
OED deprioritizes.

\begin{table*}[tb]
\centering
\caption{Per-state-class mean fidelity ($\pm 1\sigma$) across $20$~trials.
All strategies: $M = 137$ operators, $1{,}000$ shots, $n = 4$,
depol~$= 0.03$, readout~$= 0.02$.  ``Haar'' and ``Mixed'' columns
average over $10$~random states each.}
\label{tab:sota-full}
\begin{tabular}{lcccccc}
\toprule
Strategy & Product & W & GHZ & Haar (mean) & Mixed (mean) & Overall \\
\midrule
$K^*$          & $0.991 \pm 0.001$ & $\mathbf{0.932 \pm 0.012}$ & $\mathbf{0.502 \pm 0.017}$ & $0.397 \pm 0.040$ & $\mathbf{0.513 \pm 0.040}$ & $\mathbf{0.501}$ \\
D-optimal      & $0.987 \pm 0.002$ & $0.811 \pm 0.015$ & $0.225 \pm 0.020$ & $\mathbf{0.398 \pm 0.058}$ & $0.502 \pm 0.035$ & $0.479$ \\
E-optimal      & $\mathbf{0.992 \pm 0.002}$ & $0.833 \pm 0.018$ & $0.225 \pm 0.013$ & $0.317 \pm 0.034$ & $0.466 \pm 0.036$ & $0.429$ \\
A-optimal      & $\mathbf{0.993 \pm 0.002}$ & $0.833 \pm 0.017$ & $0.229 \pm 0.017$ & $0.316 \pm 0.033$ & $0.462 \pm 0.037$ & $0.428$ \\
WA-random      & $0.888 \pm 0.210$ & $0.414 \pm 0.224$ & $0.352 \pm 0.154$ & $0.242 \pm 0.149$ & $0.391 \pm 0.138$ & $0.347$ \\
DR-shadow      & $0.952 \pm 0.015$ & $0.567 \pm 0.167$ & $0.114 \pm 0.063$ & $0.066 \pm 0.017$ & $0.227 \pm 0.020$ & $0.198$ \\
Adaptive (3R)  & $0.120 \pm 0.237$ & $0.096 \pm 0.052$ & $0.139 \pm 0.201$ & $0.062 \pm 0.013$ & $0.224 \pm 0.020$ & $0.140$ \\
Adaptive (5R)  & $0.034 \pm 0.093$ & $0.092 \pm 0.115$ & $0.075 \pm 0.105$ & $0.062 \pm 0.013$ & $0.224 \pm 0.020$ & $0.133$ \\
Uniform random & $0.711 \pm 0.290$ & $0.291 \pm 0.259$ & $0.215 \pm 0.184$ & $0.063 \pm 0.019$ & $0.227 \pm 0.020$ & $0.179$ \\
\bottomrule
\end{tabular}
\end{table*}

Key observations:
\begin{itemize}
\item Product states are easy for all deterministic non-adaptive strategies
($F > 0.95$) because any structured $137$-operator subset includes enough
weight-$1$ operators; random-selection strategies (WA-random $0.888$,
uniform $0.711$) score lower due to inconsistent weight-$1$ coverage.
\item W~states show the largest $K^*$ advantage among named states:
$F = 0.932$ vs.\ $0.833$ (E-optimal) and $0.811$ (D-optimal).
$K^*$'s full saturation of weight-$1$ and weight-$2$ sectors captures the
local correlators that dominate W-state information content.
\item GHZ separates $K^*$ from all competitors: only $K^*$ maintains
enough weight-$4$ operators (16 of 81 total) to capture phase
correlations. D/E/A-optimal designs allocate weight-$4$ operators
late in the greedy sequence when their marginal contribution to
the Gram determinant/eigenvalue/trace is small.
\item Haar-random and mixed states show moderate $K^*$ advantage,
consistent with the weight-class balance
being most valuable when many weight classes carry signal.
\item WA-random's high variance reflects the sensitivity of
within-weight-class operator choice: the right weight allocation
helps, but random selection within each class produces inconsistent
results across seeds.
\end{itemize}

\subsection{Hybrid $K^*$+adaptive test}
\label{sec:hybrid-test}

A natural question is whether adaptive refinement can improve upon
$K^*$ when used as a well-informed first round.  We test this with
a hybrid protocol: round~1 measures all $137$ $K^*$ operators and
reconstructs via MLE; subsequent rounds select from the remaining
$255 - 137 = 118$ operators using the same information-gain
heuristic described above, with $2$ or $4$ adaptive rounds.

Table~\ref{tab:hybrid} compares six strategies at three budget
levels ($20$~trials, $22$~states excluding W, same noise model as
Table~\ref{tab:sota-full}).

\begin{table}[tb]
\centering
\caption{Cross-budget comparison of pure and hybrid strategies.
Overall mean fidelity across $22$~states (excluding W).
At $n = 4$: $4^n - 1 = 255$ non-identity Pauli operators.
Budget $2\times$ and $3\times$ both cap at $255$ (informationally
complete).}
\label{tab:hybrid}
\resizebox{\columnwidth}{!}{%
\begin{tabular}{lcccccc}
\toprule
Budget & $K^*$-only & $K^*{+}$Ad-2R & $K^*{+}$Ad-4R & Ad-3R & Ad-5R & Random \\
\midrule
$137$ (54\%) & $\mathbf{0.481}$ & $0.481$ & $0.481$ & $0.143$ & $0.145$ & $0.175$ \\
$255$ (100\%) & $\mathbf{0.488}$ & $0.484$ & $0.488$ & $0.189$ & $0.191$ & $0.190$ \\
\bottomrule
\end{tabular}}% end resizebox
\end{table}

Key findings:
\begin{itemize}
\item At $M = 137$ (sub-complete), the hybrid protocols
degenerate to pure $K^*$ since no budget remains for adaptive
rounds.  All three $K^*$-based strategies tie at $F \approx 0.481$.
\item At $M = 255$ (informationally complete), $K^*$-only
measures all $255$ operators in a single batch.
Hybrid $K^*{+}$Adaptive protocols measure the same $255$
operators but split them across rounds, incurring reconstruction
overhead.  $K^*$-only wins by $+0.004$ over the best hybrid
variant (within noise), and the adaptive overhead provides no
marginal benefit.
\item Pure adaptive (Adapt-3R/5R) improves from $F \approx 0.14$
at sub-complete budget to $F \approx 0.19$ at complete budget, but
remains far below $K^*$-only ($0.49$).  The round-splitting itself
degrades performance: even with $255$~total operators and a good
$K^*$ initialization, intermediate reconstructions add noise to
the selection process without improving the final reconstruction.
\item Single-batch architecture is not a limitation but a
structural advantage: it avoids the reconstruction noise that
contaminates adaptive operator selection.
\end{itemize}

Source code: \texttt{sota\_hybrid.py}.  Full per-state results are
saved in \texttt{sota\_hybrid\_n4.json}.

\subsection{Statistical significance and reconstructor dependence}
\label{sec:significance}

To assess whether the $K^*$ advantage over D-optimal is
statistically significant, we re-ran the comparison with per-trial
paired fidelities ($20$~trials, same noise model) and applied the
Wilcoxon signed-rank test (one-sided, $H_1\colon F_{K^*} > F_\text{D-opt}$).
We additionally tested three reconstructors: robust MLE, linear
inversion (project-and-truncate), and ordinary least-squares.

\begin{table}[tb]
\centering
\caption{Paired Wilcoxon signed-rank test: $K^*$ vs D-optimal
under three reconstructors.  $\Delta\bar{F}$ is the mean
per-trial fidelity difference.  All $p$-values are one-sided.}
\label{tab:significance}
\begin{tabular}{lcccl}
\toprule
Reconstructor & $\bar{F}_{K^*}$ & $\bar{F}_\text{D-opt}$ & $\Delta\bar{F}$ & $p$-value \\
\midrule
Robust MLE    & $0.502$ & $0.480$ & $+0.022$ & $< 10^{-4}$ \\
Linear inv.   & $0.454$ & $0.437$ & $+0.017$ & $< 10^{-4}$ \\
Least-squares & $0.454$ & $0.437$ & $+0.018$ & $< 10^{-4}$ \\
\bottomrule
\end{tabular}
\end{table}

$K^*$'s advantage is highly significant ($p < 10^{-4}$) under all
three reconstructors (Table~\ref{tab:significance}), confirming
that the $+0.022$ margin in Table~\ref{tab:sota-full} is not a
statistical fluctuation.  This gap is consistent across
reconstructors ($+0.017$--$+0.022$), indicating that the advantage
originates in operator selection, not in a reconstructor-specific
interaction.

\paragraph{Random selection under linear methods.}
Uniform random selection ($\bar{F} = 0.574$) \emph{outperforms}
both $K^*$ ($0.454$) and D-optimal ($0.437$) under linear
inversion and least-squares.  This inversion occurs because
linear methods do not enforce positivity during reconstruction:
they solve a linear system and project the result onto the
positive semidefinite cone.  Random selection's broader Pauli
coverage produces a better-conditioned design matrix for the
linear solve, while $K^*$'s structured weight-class allocation
creates a less uniform design that linear methods cannot exploit.

Under MLE, the iterative positivity constraint converts $K^*$'s
weight-class balance into a tighter likelihood surface: the
structured allocation constrains the feasible set of density
matrices more effectively than uniform coverage.  This explains
why $K^*$'s advantage ($+0.321$ over random under MLE) vanishes
and inverts ($-0.120$) under linear inversion.

The practical implication is that the $K^*$ prescription is designed
for, and should be used with, nonlinear reconstructors (MLE,
Bayesian mean estimation) that enforce physicality.  For linear
estimators on near-complete measurement sets, uniform random
selection suffices.

Source code: \texttt{significance\_and\_reconstructor.py}.

\subsection{Non-uniform shot allocation}
\label{sec:shot-alloc}

A potential confound is that all operators receive equal shots
($1{,}000$ each), while high-weight operators have lower SNR due
to readout error compounding as $(1 - 2\epsilon)^w$.  We test
whether weight-scaled shot allocation changes the $K^*$ advantage.

The weight-scaled allocation assigns shots proportional to
$1/(1 - 2\epsilon)^{2w}$ per operator, normalized so the total
budget remains $137{,}000$ shots.  At $\epsilon = 0.02$, this
gives: weight-$0$: $812$ shots, weight-$1$: $881$, weight-$2$:
$956$, weight-$3$: $1{,}037$, weight-$4$: $1{,}124$.  The
redistribution is modest because $\epsilon = 0.02$ produces only
a $1.38\times$ difference between weight-$0$ and weight-$4$.

Over $20$~trials on $22$~states ($n = 4$, same noise model as
Table~\ref{tab:sota-full}):
\begin{itemize}
\item $K^*$ fidelity: $0.476$ (uniform) vs.\ $0.474$
(weight-scaled); Wilcoxon $p = 0.21$ (not significant).
\item Random fidelity: $0.454$ (uniform) vs.\ $0.454$
(weight-scaled); Wilcoxon $p = 0.59$ (not significant).
\item $K^*$-vs-random gap: $+0.021$ (uniform) vs.\ $+0.020$
(weight-scaled); Wilcoxon $p = 0.43$ (not significant).
\end{itemize}

The $K^*$ advantage is driven by operator selection, not shot
allocation.  At larger readout error rates (e.g., $\epsilon = 0.05$),
weight-scaling would redistribute shots more aggressively, but the
operator selection effect would remain independent.

Source code: \texttt{shot\_allocation\_test.py}.

\subsection{Scaling to $n = 5$}
\label{sec:scaling-n5}

At $n = 5$, $K^* = 5$ and $M = N_5(5) = 333$ operators, giving
coverage $332/1023 = 32.5$\% of non-identity Paulis.  We compare
$K^*$, weight-allocated random (WA-random), and uniform random over
$10$~trials on $13$~states (product, W, GHZ, $5$~Haar-random pure,
$5$~mixed rank-$4$).

\begin{table}[tb]
\centering
\caption{$n = 5$ comparison ($M = 333$, $10$~trials, depol~$= 0.03$,
readout~$= 0.02$).}
\label{tab:n5-scaling}
\begin{tabular}{lcccccc}
\toprule
Strategy & Product & W & GHZ & Haar & Mixed & Overall \\
\midrule
$K^*$     & $0.986$ & $0.843$ & $0.497$ & $0.312$ & $0.378$ & $0.444$ \\
WA-random & $0.981$ & $0.753$ & $0.535$ & $0.307$ & $0.391$ & $0.443$ \\
Uniform   & $0.507$ & $0.355$ & $0.199$ & $0.032$ & $0.117$ & $0.139$ \\
\bottomrule
\end{tabular}
\end{table}

$K^*$ dominates uniform random by $+0.305$ in overall fidelity
(Table~\ref{tab:n5-scaling}).  The gap versus WA-random narrows to
$+0.001$ overall, indicating that weight-class allocation captures
${\sim}\,99$\% of the advantage at $n = 5$ (compared to $81$\% at
$n = 4$; main text Sec.~V\,B).  The specific
intra-class Krawtchouk selection retains a meaningful edge on
W~states ($+0.090$), where local structure rewards the particular
operators $K^*$ selects within weight-$1$ and weight-$2$.

The dimensional dependence prediction
(Sec.~\ref{sec:dim-proof}) is confirmed: as $n$ grows past the
sweet spot at $n = 4$, the weight allocation becomes increasingly
sufficient and intra-weight selection matters less.  Greedy
D-optimal was not tested at $n = 5$ due to the $O(4^n \cdot M)$
computational cost of the greedy search over $1{,}024$ candidates.

Source code: \texttt{sota\_scaling\_n5.py}.

\subsection{Caveats}

The greedy OED implementations are polynomial-time heuristics, not
globally optimal solutions (which are NP-hard).  Better OED
implementations (e.g., convex relaxation followed by rounding) might
narrow the gap with $K^*$.  The adaptive protocol uses a simple
variance-based information-gain heuristic; Bayesian-optimal adaptive
designs~\cite{Huszar2012} may perform better, particularly in regimes
closer to informational completeness.  The hybrid test
(Sec.~\ref{sec:hybrid-test}) shows that even with an optimal
$K^*$ initialization, information-gain adaptive refinement provides
no improvement; whether Bayesian-optimal adaptive designs would
fare differently at larger $n$ or with more shots remains open.
Locally-biased classical shadows~\cite{Hadfield2022} share the most
structural similarity with $K^*$: both are non-adaptive and
weight-structured.  The present comparison uses the generic
derandomized shadow protocol (greedy set cover over all $255$
target observables) rather than a locally-biased variant.  A
locally-biased shadow protocol tuned to $K^*$'s weight budget would
be an informative additional comparator; we leave this for future
work.
All comparisons use the same
robust MLE reconstructor; different estimators (compressed sensing,
Bayesian mean) might change the relative ranking.  The results are
specific to $n = 4$, $M \leq 255$, and the stated noise model.


%% ============================================================
%% Former SM S12
%% ============================================================
\section{Additional methodological details}
\label{sec:additional-details}

\subsection{MLE convergence behavior}
\label{sec:convergence}

In practice, every hardware reconstruction (both structured and random
arms, all states) reaches the $1{,}000$-iteration cap: the
$10^{-7}$ Frobenius tolerance is never met, as expected for
underdetermined systems ($137 < 255$ free parameters).  The
structured arm nonetheless converges functionally: fidelity
plateaus by iteration~${\sim}300$ with per-step changes
$\|\Delta\rho\|_F \sim 0.03$, while the random arm oscillates
at $\|\Delta\rho\|_F \sim 0.3$ without stabilizing.  Both
arms use identical code and cap; the fidelity
difference is not a convergence artifact.

\subsection{PSD contraction for constrained MLE}
\label{sec:psd-contraction}

The following result, used in the proof of
Theorem~2(iii) in the main text, bounds the
positivity-constrained MLE's squared deviation from the
unconstrained optimum relative to the true state's,
on measured coordinates.

\pclaim{env:prop:psd_contraction}{proved}{*}{prop:psd_contraction}%
\begin{proposition}[PSD contraction]
\label{prop:psd_contraction}
Let $\mathcal{S}$ be a set of Pauli operators, each measured $N_s$ times.
Let $\hat{y}_P = 2n_{P,+}/N_s - 1$ be the sample mean for
each $P \in \mathcal{S}$, let $x_P = \mathrm{tr}(P\rho)$ be the true
expectation, and let $\hat\rho$ be the PSD-constrained MLE
with Pauli coordinates $y_P(\hat\rho) = \mathrm{tr}(P\hat\rho)$.
Assume the bounded-Bloch regime
$|x_P|, |\hat{y}_P|, |y_P(\hat\rho)| \le t_\star < 1$ for all
$P \in \mathcal{S}$
(under $\mathcal{E}_\delta$, the first two bounds hold with
$t_\star = \max_P |\hat{y}_P| + \sqrt{2\ln(2M/\delta)/N_s}$,
which is computable from the sample means and guarantees
$\max_P |x_P| \le t_\star$, provided $\max_P |x_P| < 1$;
the third is a verifiable output property of the PSD-constrained MLE).
Then
\begin{equation}\label{eq:sm-psd-contraction}
  \sum_{P \in \mathcal{S}}
  \bigl(y_P(\hat\rho) - \hat{y}_P\bigr)^2
  \;\le\;
  \left(\frac{1+t_\star}{1-t_\star}\right)^{\!2}
  \sum_{P \in \mathcal{S}}
  \bigl(x_P - \hat{y}_P\bigr)^2.
\end{equation}
\end{proposition}

\begin{proof}
The per-operator log-likelihood
$\ell_P(y) = n_{P,+}\log\frac{1+y}{2}
+ n_{P,-}\log\frac{1-y}{2}$
has second derivative
$\ell_P''(y) = -n_{P,+}/(1{+}y)^2 - n_{P,-}/(1{-}y)^2$
for $|y| < 1$.
On $|y| \le t_\star$, each factor $(1 \pm y)^2$ lies in
$[(1{-}t_\star)^2,\,(1{+}t_\star)^2]$, so
$N_s/(1{+}t_\star)^2 \le |\ell_P''(y)| \le N_s/(1{-}t_\star)^2$.
Since $\hat{y}_P$ maximizes $\ell_P$
(i.e., $\ell_P'(\hat{y}_P) = 0$),
Taylor's theorem with Lagrange remainder gives,
for any $y$ with $|y|, |\hat{y}_P| \le t_\star$:
\begin{align*}
  \ell_P(y) &= \ell_P(\hat{y}_P)
    + \tfrac{1}{2}\,\ell_P''(c_P)\,(y - \hat{y}_P)^2 \\
  &\le \ell_P(\hat{y}_P) - \tfrac{N_s}{2(1{+}t_\star)^2}(y - \hat{y}_P)^2
    \quad \text{(upper bound, $|\ell_P''| \ge N_s/(1{+}t_\star)^2$),} \\
  \ell_P(y) &\ge \ell_P(\hat{y}_P)
    - \tfrac{N_s}{2(1 - t_\star)^2}(y - \hat{y}_P)^2
    \quad \text{(lower bound, $|\ell_P''| \le N_s/(1{-}t_\star)^2$),}
\end{align*}
where $c_P$ lies between $y$ and $\hat{y}_P$
(hence $|c_P| \le t_\star$).
Apply the upper bound at $y = y_P(\hat\rho)$ and the lower
bound at $y = x_P$, then sum over $P \in \mathcal{S}$.
Writing $L(\sigma) = \sum_P \ell_P(\mathrm{tr}(P\sigma))$:
\begin{align*}
  L(\hat\rho) &\le L(\hat{y})
    - \tfrac{N_s}{2(1{+}t_\star)^2}\textstyle\sum_{P \in \mathcal{S}}
    (y_P(\hat\rho) - \hat{y}_P)^2, \\
  L(\rho) &\ge L(\hat{y})
    - \tfrac{N_s}{2(1 - t_\star)^2}\textstyle\sum_{P \in \mathcal{S}}
    (x_P - \hat{y}_P)^2,
\end{align*}
Since $\rho$ is feasible ($\rho \ge 0$,
$\mathrm{tr}(\rho) = 1$), the MLE satisfies
$L(\hat\rho) \ge L(\rho)$.  Chaining the three inequalities:
\[
  L(\hat{y}) - \tfrac{N_s}{2(1{+}t_\star)^2}
  \sum_{P \in \mathcal{S}} (y_P(\hat\rho) - \hat{y}_P)^2
  \;\ge\;
  L(\hat\rho) \;\ge\; L(\rho)
  \;\ge\;
  L(\hat{y}) - \tfrac{N_s}{2(1 - t_\star)^2}
  \sum_{P \in \mathcal{S}} (x_P - \hat{y}_P)^2,
\]
which after cancelling $L(\hat{y})$, negating, and dividing by
$N_s/(2(1{+}t_\star)^2)$ yields~\eqref{eq:sm-psd-contraction}.
In the shot-noise-limited regime ($\max_P |x_P| \ll 1$),
$t_\star \to 0$ and the prefactor
$\bigl(\frac{1+t_\star}{1-t_\star}\bigr)^{2}
= 1 + O(t_\star)$ absorbs into the
higher-order error of the assembled Hilbert--Schmidt bound.

The constant factor in the reconstruction error bound
follows by the triangle inequality and $(a+b)^2 \le 2(a^2+b^2)$:
\[
  \sum(x_P - y_P(\hat\rho))^2
  \;\le\; 2\sum(x_P - \hat{y}_P)^2
       + 2\sum(\hat{y}_P - y_P(\hat\rho))^2
  \;\le\; \Bigl(2 + 2\Bigl(\frac{1+t_\star}{1-t_\star}\Bigr)^{\!2}\Bigr)
          \sum(x_P - \hat{y}_P)^2,
\]
where the second step applies~\eqref{eq:sm-psd-contraction}.
In the shot-noise-limited regime ($\max_P |x_P| \ll 1$,
so $t_\star \to 0$), the prefactor equals
$\frac{4(1+t_\star^2)}{(1-t_\star)^2} = 4 + O(t_\star)$;
the manuscript's theorem retains the exact
coefficient~$8(1{+}t_\star^2)/(1{-}t_\star)^2
= 4(1{+}t_\star^2)/(1{-}t_\star)^2 \times 2$ (Hoeffding).
The bound is additionally conservative in practice: the PSD
constraint shrinks coordinates toward zero, typically
placing $y_P(\hat\rho)$ between $\hat{y}_P$ and $x_P$.
\end{proof}

\subsection{GHZ hedge-parameter insensitivity}
\label{sec:hedge-scan}

The GHZ full-state MLE deficit ($F \approx 0.50$;
Sec.~V\,D of the main text) is structural, not an artifact of
regularization.  On hardware (IBM \texttt{ibm\_fez}), the raw
phase correlators $|\langle X^{\otimes 4}\rangle| = 0.864$ and
$|\langle Y^{\otimes 4}\rangle| = 0.864$ confirm ${\sim}\,85$\%
phase coherence is retained, ruling out decoherence.  The deficit
arises because $126$ of $137$ $K^*$ operators have near-zero GHZ
expectation, overwhelming the $11$ informative operators in the
underdetermined MLE\@.

\subsection{Identity-weight pathology at $n = 2$}
\label{sec:identity-pathology}

The variance weight $w_i \propto 1/\sigma_i^2$ diverges for the
identity operator ($b = 1$, clipped to $0.999$), which dominates
${\sim}\,92$\% of the total weight at $n = 2$ ($M = 9$), suppressing
informative correlators.  At $n = 4$ ($M = 137$), the identity
fraction drops to ${\sim}\,1$\% and does not affect reconstruction
quality.
This pathology is entirely a variance-weighting artifact: the
$R$-operator pole at $y_i = 0$ is handled by ratio clipping
(Eq.~(1) of the main text), so the issue is not
an uncontrolled divergence but rather the divergent weight assigned to
a deterministic measurement outcome.
This explains the Bell pair result in the main text: the
raw correlators are high-quality ($\langle XX\rangle = 0.942$,
$\langle ZZ\rangle = 0.970$; direct fidelity
estimate $F_{\rm direct} = 0.962$), but the identity operator
dominates the $R\rho R$ variance weight at $n = 2$, suppressing
informative correlators.  The Bell pair thus serves as a negative
control for the MLE mechanism: the identity-weight dominance predicted
by the variance-weighting formula suppresses the structured advantage
at $n = 2$, confirming that the large W-state advantage at $n = 4$
requires both weight-class completeness and sufficient dilution of
the identity term.

\subsection{Exact Hradil form for the Pauli binomial likelihood}
\label{sec:hradil-comparison}

Eq.~(1) of the main text uses the ratio-form $R$-operator
$R_{\mathrm{ratio}} \propto \sum w_i\, \mathrm{clip}(b_i/y_i,\, {-}5,\, 5)\, P_i$,
a standard iterative MLE update~\cite{Hradil1997,Paris2004}.
This form has a pole at $y_i = 0$, handled by clipping ratios
to $[-5, 5]$.
All hardware results in the main text (Table~I, Table~III) use
the ratio form with identical hyperparameters for both arms.

The exact Hradil $R$-operator, obtained by summing the POVM
contributions $(f_{i,\pm}/p_{i,\pm})\,\Pi_{i,\pm}$ for the
per-operator binomial POVM
$\Pi_{i,\pm} = (I \pm P_i)/2$, is
\begin{equation}\label{eq:hradil-exact}
  R = \zeta\, I + \frac{1}{M + \zeta\, d}
  \sum_{i=1}^{M} w_i\!\left[
  \frac{1 - b_i\, y_i}{1 - y_i^2}\, I
  + \frac{b_i - y_i}{1 - y_i^2}\, P_i\right].
\end{equation}
Both coefficients are regular at $y_i = 0$ (reducing to $I + b_i P_i$),
so no ratio clipping is required, and $R = I$ at the fixed point
$b_i = y_i$.

Both forms share the same fixed-point set (the constrained MLE),
but in the underdetermined regime ($M < d^2 - 1$) the MLE has
a manifold of solutions and the two iteration schemes act as
different implicit regularizers, converging to different points
on that manifold.
For states whose informative subspace is dense
(W, product: $\ge 35/137$ nonzero expectations), the ratio form's
aggressive updates at $y_i \approx 0$ provide effective
implicit regularization, and the ratio form is preferred for
its faster convergence.
For states whose informative subspace is sparse
(GHZ: only $11/137$ nonzero expectations), the ratio form's
clipping at the $126$ uninformative operators injects structured
noise that drives the iterate toward a near-maximally-mixed state.
The full form handles these operators correctly and recovers
a high-fidelity estimate.

On the 4-qubit hardware GHZ data (IBM \texttt{ibm\_fez},
$1{,}000$ shots, $137$ $K^*$ operators, identical raw expectations):
the ratio form gives $F = 0.496$ while the full form gives
$F = 0.887$, matching the weight-class fingerprint classifier
(Sec.~V.D of the main text) from a different direction.

Table~\ref{tab:mle-comparison} extends this comparison to every
hardware dataset with raw expectation values.
The ratio form matches or exceeds the full form on all W-state and
product datasets (well-matched informative subspace),
while the full form recovers higher fidelity only for GHZ and
Bell (sparse or identity-dominated regimes).
The Bell pair's low ratio-form fidelity ($F = 0.518$) reflects the
identity-weight pathology (Sec.~\ref{sec:identity-pathology}), not
hardware noise: its raw correlators are excellent
($\langle XX\rangle = 0.942$, direct $F = 0.962$).

\begin{table}[h]
\caption{Ratio-form vs.\ full-Hradil MLE fidelity across all
hardware datasets.  $\Delta = F_{\mathrm{Hradil}} - F_{\mathrm{ratio}}$;
negative $\Delta$ means the ratio form is preferred.
All reconstructions use identical hyperparameters
($\zeta = 0.05$, $\alpha = 0.5$).
Multi-run entries report the mean over independent runs.}
\label{tab:mle-comparison}
\begin{ruledtabular}
\begin{tabular}{lccrrr}
Dataset & $n$ & $M$ & $F_{\mathrm{ratio}}$ & $F_{\mathrm{Hradil}}$ & $\Delta$ \\
\midrule
$|{+}\rangle^{\otimes 4}$ & 4 & 137 & 0.996 & 0.926 & $-$0.070 \\
GHZ$_4$                    & 4 & 137 & 0.496 & 0.887 & $+$0.392 \\
Bell $|\Phi^+\rangle$      & 2 & 9   & 0.518 & 0.954 & $+$0.436 \\
W$_4$ (single)             & 4 & 137 & 0.891 & 0.844 & $-$0.046 \\
W$_4$ ($\times 4$ runs)    & 4 & 137 & 0.872 & 0.760 & $-$0.112 \\
W$_4$ (3-arm)              & 4 & 137 & 0.903 & 0.872 & $-$0.031 \\
W$_4$ (Rigetti)            & 4 & 137 & 0.815 & 0.681 & $-$0.135 \\
\end{tabular}
\end{ruledtabular}
\end{table}

\subsection{Dimensional dependence: linear inversion cross-check}
\label{sec:lininv-crosscheck}

Under linear inversion, the dimensional progression values are
$\Delta(S{-}AR) = 0.128 \to 0.003 \to 0.010$ for
$n = 2, 3, 4$; the $n = 2$ value remains the largest, but the
progression is non-monotonic ($n{=}3 < n{=}4$), so the smooth
decrease is MLE-specific.  That the smooth progression appears only
under MLE is itself informative: it confirms that the advantage
mechanism is nonlinear (variance weighting, positivity projection)
rather than a simple bias in the operator count.

The large drop in $\Delta(S{-}AR)$ between $n = 3$ and $n = 4$ is
explained by low-weight saturation: for $n \geq 4$ (where $K^* \geq 5$),
the $K^*$ selection includes every weight-$0$, weight-$1$, and
weight-$2$ Pauli operator, so structured and weight-allocated random
sets are identical at these weights and the residual arises entirely
from weight $\geq 3$ differences.  For $n = 2, 3$ (where $K = n < K^*$),
selection freedom exists at all weights including the informationally
dominant weight-$1$ sector.

\subsection{Measurement paradigm comparison}
\label{sec:paradigm-table}

Table~\ref{tab:regimes-sm} maps each measurement paradigm to its
optimal operating regime.

\begin{table}[tb]
\centering
\caption{Measurement paradigms and their optimal operating regimes.
``Budget'' is relative to informational completeness ($4^n - 1$).
``Interactive'' means the protocol requires mid-experiment feedback.}
\label{tab:regimes-sm}
\begin{tabular}{p{2.0cm}p{2.4cm}p{3.2cm}}
\toprule
Paradigm & Optimal regime & Practical scenario \\
\midrule
$K^*$ (this work) &
  Sub-complete budget, unknown state, single batch &
  Cloud QPU tomography; pay-per-circuit; no interactive access \\[4pt]
Classical shadows~\cite{Huang2020} &
  Few target observables, large $n$ &
  VQE/QAOA energy estimation; specific correlator measurements \\[4pt]
Derand.\ shadows~\cite{Huang2021} &
  Known target observable set, moderate $n$ &
  Hamiltonian learning; partial tomography of known subsystems \\[4pt]
Adaptive Bayesian~\cite{Huszar2012} &
  Overcomplete budget ($M \gg 4^n$), interactive access &
  Lab-based small-$n$ characterization with fast feedback loop \\[4pt]
Shadow overlap &
  Fidelity certification (pass/fail) &
  Production-line device verification; $O(n^2)$ measurements suffice \\[4pt]
DFE / witnesses &
  Known target state family &
  ``Is this a GHZ?'' ($2$--$4$ circuits, not $137$) \\[4pt]
Compressed sensing &
  Low-rank prior, large $n$ &
  Large-$n$ states with known rank structure \\
\bottomrule
\end{tabular}
\end{table}

The $K^*$ regime (single-batch, state-independent, sub-complete,
full reconstruction) is the default workflow on current cloud
platforms (IBM, Rigetti, IonQ), where circuit submissions are
non-interactive and priced per shot.  As qubit counts grow, the
ratio $M/4^n$ shrinks exponentially, making the underdetermined
regime inescapable for full-state tomography at any practical budget.

Outside this regime, purpose-built methods are preferable.
Shadow-based methods dominate when only a polynomial number of
observables are needed~\cite{Huang2020}.
Adaptive protocols dominate when interactive access is available
and per-round budgets are overcomplete~\cite{Huszar2012}: at $n = 2$
($M = 9$, $4^n - 1 = 15$), two adaptive rounds can nearly saturate
informational completeness.

\subsection{Extended limitations and robustness checks}
\label{sec:extended-limitations}

\paragraph{Intra-weight ordering robustness.}
Seven orderings with identical weight budgets show a
bimodal distribution: four orderings (including the default) cluster
at $F \approx 0.80$--$0.82$; three adversarial orderings (reverse,
antistride) that anti-correlate with support-pattern spread give
$F \approx 0.31$--$0.43$.  The default sits at the median of the
typical cluster (see Sec.~\ref{sec:ordering} for per-ordering data).

\paragraph{Parameter-free qualification.}
\label{sec:param-free}
The ``parameter-free'' claim refers to the measurement design:
the weight-level budget $\{1, 12, 54, 54, 16\}$ is determined by
spectral geometry with no tuning.  The MLE reconstructor uses
standard hyperparameters (hedge $\zeta = 0.05$, damping $= 0.5$)
held constant across all arms, state classes,
and hardware platforms.
A full hyperparameter sweep varies each axis independently:
$7$~hedge values $\times$ $6$~states $= 42$ conditions
(damping fixed at $0.5$), plus $6$~damping values $\times$
$6$~states $= 36$ conditions (hedge fixed at $0.05$), for
$78$~total conditions
($\zeta \in \{0, 0.001, 0.01, 0.05, 0.1, 0.2, 0.5\}$;
damping $\in \{0.1, 0.3, 0.5, 0.7, 0.9, 1.0\}$;
six states including three Haar-random; $10$~trials
each at $p = 0.03$, readout $2$\%, $10^3$~shots).
The $K^*$ advantage sign is positive for W, GHZ, and product
states across every setting tested.
F($K^*$, W) ranges from $0.72$ (damping~$= 0.1$) to $0.86$
(damping~$= 0.7$); the default ($\zeta = 0.05$, damping~$= 0.5$)
sits near the center of this range, not at an optimized extreme.

\paragraph{M=255 hardware control.}
An informationally complete hardware control ($M = 255$, all
nontrivial Paulis) was not run; simulation at $M = 255$ gives
$F = 0.488$ overall across $22$~states
(Sec.~\ref{sec:hybrid-test}, Table~\ref{tab:hybrid}),
only $+0.007$ above $K^*$ at $M = 137$.  The marginal return
from the additional $118$~circuits ($1.86\times$ cost) is
negligible, consistent with reconstruction quality being
noise-limited rather than coverage-limited at this budget.

\paragraph{Statistical power of four-run sample.}
The four-run W-state sample ($n_{\rm runs} = 4$) provides a
conservative uncertainty estimate; with $3$~degrees of freedom,
a paired $t$-test gives $t(3) = 9.3$, $p = 0.003$
($95$\% CI on $\Delta F$: $[0.22, 0.44]$), excluding zero by
a wide margin.  The distribution-free sign test gives
$p = 0.0625$ (all four positive); the $95$\% confidence
interval on $\sigma$ is $[0.011, 0.066]$.  The GHZ and product results
are single-seed; the GHZ result ($\Delta F = -0.002$) is predicted
by eigenvalue-mass analysis and is insensitive to seed; the
product result ($\Delta F = +0.016$) has limited room for variation
given both arms exceed $F = 0.98$.

\paragraph{Shot-noise decomposition.}
The $K^*$ arm's run-to-run standard deviation
($\sigma_{K^*} = 0.021$) is comparable to its within-run
shot-noise bootstrap standard deviation
($\sigma_{K^*}^{\mathrm{boot}} \approx 0.019$),
indicating that $K^*$ reconstructions are stable across
independent hardware sessions; the four runs are essentially % vague-verb-justified: quantified; boot sigma=0.019 across four independent sessions
resamples of the same measurement strategy.  In contrast, the
random arm's run-to-run dispersion
($\sigma_{\mathrm{rand}} = 0.076$) exceeds its per-run
shot noise ($\sigma_{\mathrm{rand}}^{\mathrm{boot}} \approx 0.062$),
and the excess variability beyond shot noise is ${\sim}\,5$
times larger than the $K^*$ arm's excess
($\sqrt{0.076^2 - 0.062^2} \approx 0.044$ vs.\
$\sqrt{0.021^2 - 0.019^2} \approx 0.009$),
reflecting genuine variability from
different random-basis draws each run.  The paired $t$-test
remains valid because pairing absorbs the correlated noise
structure, and the DeltaF variance is dominated by the
random-arm variability, which is legitimate between-run variation.

\paragraph{Bayesian reconstruction.}
Bayesian mean estimation is a natural alternative to MLE for
underdetermined problems; we do not test it here but note that the
measurement-selection comparison is reconstructor-controlled
(same algorithm, same hyperparameters, both arms), so the
$K^*$-vs-random gap is valid within the MLE framework regardless of
whether BME would yield higher absolute fidelities.

\paragraph{Non-uniform shot allocation.}
Weight-scaled allocation at $\epsilon = 0.02$ redistributes shots
modestly ($812$--$1{,}124$ vs.\ uniform $1{,}000$) and leaves
the $K^*$-vs-random gap unchanged ($p = 0.43$, Wilcoxon;
Sec.~\ref{sec:shot-alloc}).

\paragraph{Cramér--Rao bound.}
No state-independent information-theoretic optimality bound exists for
the underdetermined tomography regime ($M < 4^n - 1$): the
Cramér--Rao bound is state-dependent and requires informational
completeness.  The hardware results
(Sec.~\ref{sec:optimality-test}) provide a direct empirical proxy:
$K^*$ outperforms all tested alternatives across two vendors, four
qubit counts, and twelve independent runs.



%% ============================================================
%% OPERATING REGIME AND LONGEVITY
%% ============================================================
\section{Operating regime and longevity analysis}
\label{sm:longevity}

This section quantifies the noise regime in which $K^*$ dominates
random Pauli sampling on raw fidelity, the regime in which the two
become comparable, and the regime in which random or shadow-based
estimators take over.  It then enumerates the structural
properties of $K^*$ that persist independently of noise level,
mapping each to specific application classes that are projected
to remain in $K^*$'s operating envelope as quantum hardware matures.

\subsection{Quantitative crossover}
\label{sm:longevity-crossover}

Define the effective per-circuit error
$\varepsilon_{\mathrm{eff}} = \varepsilon_{2Q} \cdot D$, where
$\varepsilon_{2Q}$ is the average two-qubit gate error and $D$ is
the transpiled circuit depth.
Empirically, $K^*$'s raw-fidelity advantage over a matched-budget
random Pauli set tracks $\varepsilon_{\mathrm{eff}}$:
\begin{itemize}
\item $\varepsilon_{\mathrm{eff}} \gtrsim 10^{-3}$: $K^*$ dominates
  by $\Delta F \gtrsim 0.05$.  This regime covers all
  superconducting QPUs shipping in 2026 at any depth $D \gtrsim 30$.
\item $10^{-4} \lesssim \varepsilon_{\mathrm{eff}} \lesssim 10^{-3}$:
  transition region; outcome depends on the target state's
  weight-class profile.
\item $\varepsilon_{\mathrm{eff}} \lesssim 10^{-4}$: random or
  shadow-based estimators may match or exceed $K^*$ on raw fidelity
  for generic states; the reconstruction is information-limited
  rather than noise-limited.  This is the post-fault-tolerance
  logical-qubit regime.
\end{itemize}

Table~\ref{tab:noise-regime} maps current platforms to
$\varepsilon_{\mathrm{eff}}$ at representative depths.  Every
near-term experimental workload of interest (variational chemistry,
Hamiltonian dynamics, error-mitigation pipelines, mid-circuit
measurement) sits in the $K^*$-dominant regime.

\begin{table}[h]
\centering
\caption{Effective per-circuit error
$\varepsilon_{\mathrm{eff}} = \varepsilon_{2Q} \cdot D$ for
representative platforms at three depths.  Values $\gtrsim 10^{-3}$
correspond to the $K^*$-dominant regime.}
\label{tab:noise-regime}
\begin{tabular}{lcccc}
\hline\hline
Platform & $\varepsilon_{2Q}$ & $D{=}30$ & $D{=}100$ & $D{=}300$ \\
\hline
IBM Heron (2026)        & $0.005$  & $0.15$  & $0.50$  & $1.50$  \\
Rigetti Ankaa-3 (2026)  & $0.015$  & $0.45$  & $1.50$  & $4.50$  \\
Quantinuum H2 (2026)    & $0.001$  & $0.03$  & $0.10$  & $0.30$  \\
IonQ Forte (2026)       & $0.004$  & $0.12$  & $0.40$  & $1.20$  \\
Projected FT (2030+)    & $10^{-6}$ & $3{\times}10^{-5}$
                                    & $10^{-4}$
                                    & $3{\times}10^{-4}$ \\
\hline\hline
\end{tabular}
\end{table}

\subsection{Use-case taxonomy}
\label{sm:longevity-usecases}

The noise-regime crossover induces a use-case split that is
expected to remain stable as the noise floor falls.

\paragraph*{$K^*$ dominates (durable):}
\begin{itemize}
\item Variational and dynamics workloads at depth $\gtrsim 50$
  (NISQ chemistry, Hamiltonian simulation, QAOA at moderate $p$).
\item Error-mitigation regimes (zero-noise extrapolation,
  probabilistic error cancellation, virtual distillation), which
  amplify shot noise to suppress gate noise and effectively push
  any platform back into the noise-dominated regime.
\item Mid-circuit measurement and feedback experiments, which
  introduce classical-control noise channels not captured by
  $\varepsilon_{2Q}$.
\item Compositional reconstruction at any $n \ge 6$.  The
  $18\times$ operator-count saving at $n = 8$
  (main text, Limitations subsection) grows monotonically
  with $n$ and is independent of per-gate fidelity.
\item Hamiltonian learning and time-crystal certification, where
  off-diagonal correlators (XY, partial-transpose negativity,
  $\ell_1$ coherence) are the observable of interest rather than
  state fidelity.
\item Qudit and $q \ge 3$ hardware, where the Krawtchouk bridge
  (Sec.~\ref{sec:qudit}) gives $K^*$ a structural derivation that
  random or shadow-based protocols lack.
\item Certifiable benchmarks for procurement, regulatory, or
  contested deployments where worst-case bounds (not average-case
  efficiency) are required.
\item Code-block characterization at the surface-code patch scale.
  The single-patch reconstruction problem stays at $n = 4$--$9$
  physical qubits (the regime $K^*$ was constructed for) for the
  lifetime of every superconducting and photonic FT platform.
\end{itemize}

\paragraph*{Random or shadow-based methods take over:}
\begin{itemize}
\item Logical-qubit characterization post-fault-tolerance
  (logical error rates $\sim 10^{-6}$ per gate).
\item Trapped-ion shallow circuits ($D \le 5$, $n \le 6$) on
  current Quantinuum H-class hardware, where shot noise already
  dominates gate noise.
\item Photonic source benchmarking with near-deterministic
  single-photon sources at very small $n$.
\item State-preparation verification of near-product states
  ($F \gtrsim 0.99$), where shot-noise efficiency dominates
  reconstruction quality.
\end{itemize}

\subsection{Why the structural advantages do not depreciate}
\label{sm:longevity-structural}

Three properties of $K^*$ are noise-independent and therefore
unaffected by improvements in gate fidelity.

\paragraph*{(i) Certifiable worst-case bounds.}
The basin separation theorem (Theorem~2 of the main text) provides
a \emph{worst-case} fidelity guarantee on the informative
subspace.  Random Pauli ensembles and classical-shadow estimators
provide \emph{expectation-value} efficiency bounds, which are not
interchangeable with worst-case bounds.  No amount of additional
samples converts an average-case shadow estimator into a worst-case
certified one.  Applications that require certifiable bounds
(safety-critical control, regulatory compliance, contested
benchmarks) therefore have no alternative to $K^*$-class methods.

\paragraph*{(ii) Single-shot off-diagonal access.}
$K^*$ includes weight-$1$ and weight-$2$ Pauli operators across
all single-qubit Pauli types (not only $Z$), so off-diagonal
correlators such as
$\langle X_i X_j \rangle$, $\langle X_i Y_j \rangle$, and
$\langle Y_i Y_j \rangle$ are available as raw expectation values
rather than as post-processed reconstructions.
The Floquet DTC application (Sec.~V\,H of the main text)
illustrates this with Pearson $r = 0.94$ between hardware and
ideal correlators.
Shadow-based estimators can produce the same correlators only
through post-processing whose variance scales with the inverse
median overlap between the random measurement basis and the target
operator, a penalty that is independent of gate noise and that
$K^*$ avoids by construction.

\paragraph*{(iii) Compositional scale invariance.}
$K^*$'s patch-based architecture uses a constant operator budget
per 4-qubit patch ($137$ operators), and the number of patches
scales linearly with the system size and topology, not with the
qubit count.  Random Pauli ensembles need $O(4^n)$ samples in the
worst case to maintain coverage; classical shadows need
$O(\log d)$ samples for linear functionals but no constant-cost
recipe for full local density matrices.  The compositional gap
($1\times$ at $n = 4$, $4.9\times$ at $n = 6$, $18\times$ at
$n = 8$) grows monotonically and is independent of per-gate
fidelity.

\subsection{Transition through fault tolerance}
\label{sm:longevity-ft}

Fault-tolerant quantum computation does not eliminate the physical
tomography problem; it relocates it.
Surface-code patches at distance $d = 3, 5, 7$ comprise
$9, 25, 49$ physical qubits, but the \emph{single-patch}
characterization problem (calibration, syndrome decoding research,
post-mortem analysis of code-block failures) operates on $4$--$9$
physical qubits at any given time, exactly the regime where
$K^*$ was constructed.
The protocol therefore transitions from a near-term performance
tool to a long-term verification primitive embedded in the
fault-tolerant stack, rather than being superseded.
The mathematical results (Theorems~1 and~2 of the main text and
the Krawtchouk--measurement bridge of
Sec.~\ref{sec:qudit}) are theorems and do not depreciate; they
become standard background in coding-theoretic quantum information.

%% ============================================================
%% OPEN QUESTIONS
%% ============================================================
\section{Open questions}
\label{sm:open}

\begin{enumerate}
\item \textbf{MLE with circuit-compressed measurements.}
Does the MLE advantage ($F \approx 0.87$ vs.\ $0.36$ for linear
inversion) persist when using the 29-circuit compressed measurement
set with 196 operators?  This would combine circuit efficiency with
nonlinear reconstruction quality.

\item \textbf{Qudit hardware validation.}
The two-qutrit case ($q = 3$, $n = 2$: 29 operators out of 81) is
ready for experimental test on trapped-ion platforms with native
qutrit support.  Simulation predicts $1$-RDM fidelity
$F = 0.972 \pm 0.028$ for the W~state
(Sec.~\ref{sec:qudit}); hardware confirmation would
validate the $q$-ary Krawtchouk extension beyond simulation.

\item \textbf{Optimal $K$ for qudits.}
Is there a principled criterion (analogous to weight-1 dominance
for qubits) that determines the right $K$ for general $q$?  The
current $K^* = K_{\mathrm{mass}} = q^2$ is spectral-geometric;
full weight saturation requires $K_{\mathrm{sat}} > q^2$
(e.g., $K_{\mathrm{sat}} = 14$ for $q = 3$, $n = 4$).
A closed-form expression for $K_{\mathrm{sat}}(q,n)$ and a
unified condition bridging the spectral and operational
thresholds would strengthen the framework.

\item \textbf{Scaling to larger $n$.}
The qubit results at $n = 4$ are encouraging.  Does the circuit
compression ratio improve or degrade at $n = 6, 8$?  The
tensor-product basis structure of larger systems could enable even
better compression, but the coverage ratio
$N_n(K^*) / (4^n - 1)$ shrinks exponentially.

\item \textbf{State-dependent MLE allocation.}
Under MLE reconstruction, does non-uniform shot allocation (more
shots to high-weight operators, informed by a preliminary low-shot
estimate) improve full-state fidelity?  The weight-class ablation
(the ablation table in Appendix~F of the main text) shows that removing
weight-4 operators improves full-state MLE fidelity, while
subspace estimation requires exactly those operators for GHZ
recovery ($F = 0.926$; main text Sec.~V\,D).
The weight-class fingerprint classifier introduced in
Sec.~V\,D of the main text resolves this tension: it automatically
detects GHZ-type states via the parity ratio $r = E_{\mathrm{odd}} /
E_{\mathrm{even}}$ and switches to subspace estimation, recovering
$F = 0.887$ from the same data.  Whether adaptive shot allocation
within a single measurement campaign can further improve full-state
MLE remains open.
\end{enumerate}

%% ============================================================
%% LEAN4 ADVERSARIAL VETTING
%% ============================================================
\section{Lean~4 adversarial vetting of the trust base}
\label{sm:lean-adversarial}

Axioms A1--A4 listed in the manuscript (Appendix~A) are the
user-facing surface of a broader \texttt{axiom} layer that
combines opaque type/\allowbreak{}function/\allowbreak{}predicate
interfaces (for objects that cannot be constructed from the rational
Pauli algebra) with propositional axioms encoding their classical
properties.  Every \texttt{axiom} entry is the result of an explicit
load-bearing audit: each candidate is commented out and
\texttt{lake build} re-run; an axiom is retained only if a downstream
theorem actually invokes it.  The audit covers six aspects:

\begin{enumerate}\setlength\itemsep{0pt}\tolerance=9999\emergencystretch=2em
\item[\textbf{(i)}]
\emph{Load-bearing.}
\texttt{mutate\_axioms.py}: every retained \texttt{axiom} is
commented out one at a time, full \texttt{lake build} re-run.
Every propositional axiom was killed by at least one downstream
theorem (opaque-type declarations are excluded, since removing a
type axiom triggers a trivial elaboration cascade unrelated to
proof-chain dependence). Textbook results are instantiated as
LHS-pinned axioms whose left-hand sides are opaque-typed.
\item[\textbf{(ii)}]
\emph{Hypothesis tightening.}
Lean's \texttt{linter.\allowbreak{}unusedVariables} was harvested across all
theorems; every unused propositional precondition was
removed. The resulting theorems hold for arbitrary rationals,
showing that the PSD-contraction / triangle / Hoeffding
assembly is sign-agnostic.
\item[\textbf{(iii)}]
\emph{Generalization stress.}
\texttt{generalization\_\allowbreak{}stress.py}: parses every theorem
docstring claiming parameter abstraction
and verifies the
signature exposes the corresponding free
$\mathbb N$-binder.
All claimed abstractions are confirmed present in the theorem
signatures.
\item[\textbf{(iv)}]
\emph{Definitional anchoring.}
\texttt{mutate\_defs.py}: every canonical numerical
\texttt{def} (Krawtchouk coefficients,
$r_4$-shell counts, Gram eigenvalue formula,
Hamming-distance bit slot, rank/eigenvalue
witness table) was perturbed by a numerically distinct
single edit. Every perturbation was caught by a downstream theorem,
confirming that no opaque constant can be silently rewired without
breaking the proof chain.
\item[\textbf{(v)}]
\emph{Import minimization.}
\texttt{mutate\_imports.py}: each
\texttt{import KstarFormal.X} line is commented out one
at a time and the project rebuilt, so that only imports whose
removal breaks the build are retained.
\item[\textbf{(vi)}]
\emph{Universe-polymorphism stress.}
\texttt{KstarFormal/\allowbreak{}UniverseCheck.lean}: a Lean meta-program
that walks \texttt{env.constants}, filters by module of origin,
excludes auto-generated suffixes (\texttt{.match\_*},
\texttt{.rec}, \texttt{.casesOn}, etc.), and reports any
constant whose \texttt{levelParams} is non-empty. All
user-written constants are universe-monomorphic
(\texttt{Type\,0} only).
\end{enumerate}
\begin{sloppypar}\raggedright
In addition, the algebraic deep theorems were stress-tested by a
mutation-based test suite
(\texttt{mutate\_and\_test.py}): every \texttt{deep\_*} theorem was
subjected to four classes of decidable-falsifying mutations
(literal flips, $\le$/$<$ swaps, RHS perturbations, replacement
with \texttt{False}); no tautologies were found---every theorem
was killed by at least one mutator. All audit tooling
(\texttt{mutate\_axioms.py},\allowbreak{} \texttt{mutate\_defs.py},\allowbreak{}
\texttt{mutate\_imports.py},\allowbreak{} \texttt{mutate\_and\_test.py},\allowbreak{}
\texttt{generalization\_stress.py},\allowbreak{}
\texttt{unused\_hypothesis\_detector.py}, and
\texttt{UniverseCheck.lean}) ships with the
verification repository~\cite{github_verification}.
\end{sloppypar}

\subsection{Universal-$n$ airtightness}
\label{sm:lean-universal}

\begin{sloppypar}\raggedright
The manuscript's universal-$n$ claims are formalized in Lean
at arbitrary $n$ via the following modules of
\texttt{KstarFormal}:
\texttt{Combinatorics/\allowbreak{}LatticeCountGeneric.lean}
(generic shell count $r_n(k)$, ball count $N_n(K)$, parity-weight
count $c_w(K,n)$, and Krawtchouk eigenvalue formula at arbitrary~$n$
and~$K$, with $n=4$ compatibility proved via \texttt{native\_decide});
\texttt{Universal.lean} (universal-$n$ wrappers for Lemma~1,
Lemma~2, Corollary~1, and Theorem~1 parts~(i)--(iii), plus a
universal-$K$ monotonicity proof for Lemma~3 by induction on~$n$);
\texttt{UniversalSpectralChar.lean} (universal-$n$ witness
indistinguishability and fidelity chain for Theorem~2 parts~(i)--(iii));
\texttt{UniversalKstar.lean} ($K^*{=}5$ saturation witnesses at
$n\in\{4,5,6,7,8\}$ for part~(iv), plus the full-saturation dichotomy
at $n\le5$ vs.\ $n\ge6$);
\texttt{Combinatorics/\allowbreak{}CwClosedForm.lean} (universal-$n$
symbolic closed forms $c_0(5,n)=1+2n$, $c_1(5,n)=2n(2n{-}1)$,
$c_2(5,n)=2n(n{-}1)$ proved by structural induction on~$n$);
and \texttt{Combinatorics/\allowbreak{}PatchDecomp.lean} plus
\texttt{UniversalKstarCompositional.lean} (for every $n\ge4$,
every weight-$1$ and weight-$2$ Pauli support embeds in a
4-qubit patch, delivering the compositional $K^*{=}5$ saturation
claim at universal~$n$).
\end{sloppypar}

\medskip
\noindent
Per-theorem axiom footprint (verified via \texttt{\#print axioms}):

\begin{table}[h!]
\centering
\begin{tabular}{lll}
\toprule
\textbf{Theorem} & \textbf{Lean name} & \textbf{Propositional axioms} \\
\midrule
Lemma~1  & \texttt{lem1\_hessian\_universal}
         & none (Lean core only) \\
Lemma~2  & \texttt{lem2\_purity\_universal}
         & none \\
Lemma~3  & \texttt{lem3\_eigenvalue\_monotonicity\_universal}
         & none \\
Lemma~4  & \texttt{lem4\_hypergeometric\_bound}
         & none \\
Cor.~1   & \texttt{cor1\_approx\_locality\_universal}
         & none \\
Thm.~1(i--iii) & \texttt{thm1\_basin\_universal}
         & none \\
Thm.~2(i--iii) & \texttt{thm2\_spectral\_char\_universal}
         & A2 (Uhlmann 1976), A3 (Fuchs--van~de~Graaf 1999) \\
Thm.~2(iv)     & \texttt{kstar\_saturation\_universal}
         & none (\texttt{native\_decide} at $n\in\{4,\ldots,8\}$) \\
Thm.~2(iv.c)   & \texttt{compositional\_saturation\_w12}
         & none (Lean core only, universal in $n\ge4$) \\
Thm.~3   & \texttt{thm3\_asymptotic\_separation}
         & none \\
\bottomrule
\end{tabular}
\end{table}

\noindent
All entries above are sorry-free; each \texttt{\#print axioms}
output includes only the listed propositional axiom(s) plus the
three Lean core axioms \texttt{propext}, \texttt{Classical.choice},
\texttt{Quot.sound}, and (for quantum-state theorems) the opaque
type/function declarations \texttt{DensityMatrix}, \texttt{qfidelity},
\texttt{witnessPlus}, \texttt{witnessMinus}.
Theorem~2 part~(iv.c), the compositional patch architecture for
$n\ge6$, is certified at universal~$n$ by
\texttt{compositional\_saturation\_w12}: for every $n\ge4$, every
weight-$1$ and weight-$2$ Pauli support embeds in a 4-qubit patch,
so composing the per-patch $K^*{=}5$ saturation (already verified
at $n=4$) covers both weight classes at arbitrary~$n$.

\paragraph{Machine-generated companion table.}
The table above is the narrative-friendly summary.  The complete,
machine-generated per-claim axiom footprint, regenerated on every
Lean build via \texttt{lean4/\allowbreak{}scripts/\allowbreak{}generate\_all.py},
is reproduced below.  Each row is parsed directly from
\texttt{\#print axioms} output and classified into Lean core axioms,
named propositional axioms (legend in the comment block), concrete-$n$
\texttt{native\_decide} certificates (\textsc{nd}\textsubscript{$N$} for $N$
certs), and opaque declarations (\textsc{opq}\textsubscript{$N$}).

\medskip
\begin{center}
\footnotesize
% Auto-generated by lean4/scripts/generate_sm_axiom_table.py
% DO NOT EDIT. Regenerate after any Lean change.
% Legend: \textsc{fvdg}=Fuchs--van de Graaf~1999; \textsc{u}=Uhlmann~1976.
% \nativedec{N}=N concrete-n native\_decide certs.
% \opaquedecl{N}=N opaque declarations (DensityMatrix/qfidelity-family).
% \providecommand lets this snippet be \input'd twice (e.g., main + appendix) without 'already defined' errors.
\providecommand{\nativedec}[1]{\textsc{nd}\textsubscript{#1}}
\providecommand{\opaquedecl}[1]{\textsc{opq}\textsubscript{#1}}
\begin{tabular}{@{}lll@{}}
\toprule
Claim ID & Lean theorem & Axiom footprint \\
\midrule
\texttt{lem:hessian} & \texttt{lem1\_hessian\_universal} & Lean core only \\
\texttt{thm:basin} & \texttt{thm1\_basin\_universal} & Lean core only \\
\texttt{cor:approx\_local} & \texttt{cor1\_approx\_locality\_universal} & Lean core only \\
\texttt{lem:purity\_main} & \texttt{lem2\_purity\_universal} & Lean core only \\
\texttt{lem:monotone} & \texttt{lem3\_eigenvalue\_monotonicity\_universal} & Lean core only \\
\texttt{thm:spectral-char} & \texttt{thm2\_spectral\_char\_universal} & Lean core + \textsc{fvdg}, \textsc{u}; \opaquedecl{4} \\
\texttt{lem:coupon} & \texttt{lem4\_hypergeometric\_bound} & Lean core only \\
\texttt{cor:lower\_bound} & \texttt{cor2\_operator\_lower\_bound\_n4} & Lean core + \nativedec{2} \\
\texttt{thm:asymptotic} & \texttt{thm3\_asymptotic\_separation} & Lean core only \\
\texttt{lem:spectral\_q\_main} & \texttt{lem5\_spectral\_decomposition\_n4} & Lean core + \nativedec{1} \\
\texttt{lem:completeness} & \texttt{lem6\_support\_completeness\_n4} & Lean core + \nativedec{5} \\
\texttt{prop:spectral\_q} & \texttt{bose\_mesner\_combinatorial\_iff} & Lean core only \\
\texttt{cor:support} & \texttt{bose\_mesner\_combinatorial\_iff} & Lean core only \\
\texttt{prop:spec-complete} & \texttt{lem6\_support\_completeness\_n4} & Lean core + \nativedec{5} \\
\texttt{prop:mass-threshold} & \texttt{weight12\_saturation\_K5\_all\_n} & Lean core only \\
\texttt{prop:psd\_contraction} & \texttt{psd\_contraction} & Lean core only \\
\texttt{fact:kstar\_count} & \texttt{M\_w\_K5\_sum} & Lean core + \nativedec{1} \\
\texttt{fact:kstar\_python\_lean\_bridge} & \texttt{physical\_kstar\_count\_n4} & Lean core only \\
\texttt{fact:delsarte\_dual\_K5} & \texttt{delsarteDual\_K5\_all\_pos} & Lean core + \nativedec{5} \\
\texttt{fact:kfull\_scaling} & \texttt{kfull\_n5} & Lean core + \nativedec{7} \\
\texttt{prop:hradil\_R\_operator} & \texttt{hradil\_R\_closed\_form} & Lean core only \\
\texttt{fact:weight3\_undersaturation\_K5} & \texttt{weight\_3\_not\_saturated} & Lean core + \nativedec{1} \\
\bottomrule
\end{tabular}

\end{center}

%% ============================================================
%% DATA AND CODE
%% ============================================================
\section{Data and code availability}
\label{sm:data-code}

The following materials are archived at Ref.~\cite{zenodo}:
\begin{itemize}
\item Raw hardware measurement counts (JSON format) from all
IBM Quantum and Rigetti experiments.
\item The complete 137-operator $K^*$ measurement set and the
29 tensor-product bases (machine-readable JSON).
\item Reconstruction code (\texttt{scripts/core.py},
\texttt{scripts/robust\_mle.py}) and SOTA comparison scripts
(\texttt{scripts/sota\_comparison.py}, \texttt{scripts/sota\_hybrid.py}).
\item Independent verification suite:
\texttt{scripts/independent-verification/} (claim-level checks)
and \texttt{cross-verification/} (exact and numerical proof checks).
\item Figure generation scripts (\texttt{figures/}).
\item Formal verification materials (\texttt{verification/}):
  claim registry mapping all 11 formal statements to manuscript lines
  (\texttt{proofs\_registry.yaml});
  SageMath/FLINT exact-arithmetic certified values
  (\texttt{certified\_values.json});
  and full eight-tier pipeline results (\texttt{verification\_results.json}).
  The Lean~4 source (Mathlib~v4.29.0-rc8, 0~\texttt{sorry})
  is maintained with CI enforcement at~\cite{github_verification}.
\item Jupyter notebook (\texttt{notebook/k\_star\_demo.ipynb})
  reproducing every hardware Table~I row, the Floquet DTC
  XY-correlator comparison, the seven-strategy SOTA ranking, the
  dimensional-dependence progression, and theorem-level
  numerical witnesses (Lemmas~1--6, Theorem~1 parts i/ii/iii plus
  finite-shot tail, Theorem~2, Theorem~3, Corollaries~1 and~2) in
  a single pass; runnable in-browser via Binder with no local
  install.
  The \texttt{axiom} trust base has been audited on six aspects
  (Sec.~\ref{sm:lean-adversarial}); see
  \texttt{scripts/independent-verification/} for the
  audit tooling.
\end{itemize}

\bibliography{refs}

\end{document}
